\pdfoutput=1
\documentclass[11pt]{amsart}

\usepackage{amsmath,amssymb,amsthm}
\usepackage[margin=1.25in]{geometry}
\usepackage{booktabs}
\usepackage{array}
\usepackage[hidelinks]{hyperref}

\setlength{\emergencystretch}{3em}
\renewcommand{\theenumi}{\roman{enumi}}
\renewcommand{\labelenumi}{(\theenumi)}
\hypersetup{pdftitle={Two-fold Langford sequences exist exactly when 3l >= 2d - 1}}

\theoremstyle{plain}
\newtheorem{theorem}{Theorem}[section]
\newtheorem{lemma}[theorem]{Lemma}
\newtheorem{proposition}[theorem]{Proposition}
\newtheorem{corollary}[theorem]{Corollary}
\newtheorem{conjecture}[theorem]{Conjecture}
\newtheorem*{thmA}{Theorem A}
\newtheorem*{thmB}{Theorem B}
\theoremstyle{definition}
\newtheorem{definition}[theorem]{Definition}
\newtheorem{example}[theorem]{Example}
\theoremstyle{remark}
\newtheorem{remark}[theorem]{Remark}

\newcommand{\Z}{\mathbb{Z}}
\newcommand{\SP}{\mathrm{SP}}
\newcommand{\lean}[1]{\par\smallskip
  {\small\raggedright\noindent\textit{Lean:} \texttt{#1}.\par}\smallskip}
\newcommand{\status}[1]{\par\smallskip
  {\small\raggedright\noindent\textit{Status:} #1\par}\smallskip}
\newcommand{\famtable}[8]{%
  \par\bigskip\noindent\textbf{#1.} #2; domain #4.
  Target order $(#3)$. Lean: #5.\par\nopagebreak\smallskip
  {\footnotesize\noindent
  \begin{tabular}{@{}l l l l >{\raggedright\arraybackslash}p{2.45in}@{}}
  \toprule
  block & $n_i$ & $M_i$ & type & class runs\\
  \midrule
  #6
  \bottomrule
  \end{tabular}\par}
  \nopagebreak\smallskip\noindent{\small Odd chain: #7. Even chain: #8.}\par}

\author{JD Jones}

\title[Two-fold Langford sequences]{Two-fold Langford sequences exist exactly when
$3l \ge 2d-1$, and what is known for every multiplicity}

\date{27 September 2026}

\subjclass[2020]{05A05 (Primary); 11B83, 68V20 (Secondary)}
\keywords{Langford sequences, Skolem-type sequences, $m$-fold sequences, signed
permutations, block reversals, Farkas certificates, formal verification}

\begin{document}

\begin{abstract}
A two-fold Langford sequence of order $l$ and defect $d$ is a sequence of length $4l$
in which every $p \in [d,d+l-1]$ occupies two disjoint pairs of positions at distance
$p$. Alkasasbeh, Dyer and Howell asked for necessary and sufficient conditions for the
existence of $m$-fold Langford sequences, $m \ge 2$. We answer the case $m = 2$: for
$d, l \ge 1$ a two-fold Langford sequence of order $l$ and defect $d$ exists if and only
if $3l \ge 2d-1$. Necessity is a distance count. Sufficiency reduces, by concatenation
and a mirror-symmetric ansatz, to a problem about permutations of $\{1,\dots,l\}$ whose
signed displacements are prescribed; an inversion symmetry and a composition lemma
reduce that problem to a band near its diagonal, which is closed by twenty-one
block-reversal families and the two-block reversal $\tau_l$, each proved by a class
table, together with eleven explicit permutations. For every $m$ we prove the counting bound $(2m-1)l \ge 2d-1$ with a
parity condition for odd $m$, show that the bound is attained for every $m$ and only by
sequences in which every pair straddles the midpoint, and settle the order $l = 1$ (a
sequence exists if and only if $d \mid m$). We prove a forced-endpoint bound and a family
of necessary conditions, $r(T-r) \le m\sum_{p=d}^{d+l-1}|p-T|$ with $r = ml \bmod T$ for
every integer $T \ge 1$, which contains the counting bound and the necessity half of the
order-one result. For every $m \ge 3$ the cell $(d,l) = (3m-3,3)$ satisfies counting and
parity and has no sequence, so counting and parity are not sufficient. These
statements are formalized and machine-checked in Lean~4. For $m \ge 3$ we also describe,
with exact Farkas certificates and closed-form edges, the empty bands inside the counting
bound, and give an exact finite test that detects every negative member of that
certificate family; that part is proved on paper and not formalized, and the general
characterization remains open.
\end{abstract}

\maketitle

\section{Introduction}\label{sec:intro}

Write $[u,v]$ for the set of integers $x$ with $u \le x \le v$. A \emph{Langford
sequence} of defect $d$ and order $l$ is a sequence of length $2l$ in which each
$p \in [d,d+l-1]$ occurs exactly twice, at two positions $p$ apart; at $d = 1$ these are
Skolem sequences \cite{Skolem,OKeefe}, and their existence was settled by Simpson
\cite{Simpson}. Alkasasbeh, Dyer and Howell \cite{ADH} introduced the \emph{two-fold}
version, in which every $p \in [d,d+l-1]$ occupies two disjoint pairs of positions at
distance $p$, so that the sequence has length $4l$; more generally an \emph{$m$-fold}
Langford sequence has length $2ml$ and $m$ disjoint pairs for every $p$ (the $d = 1$ case
is the $m$-fold Skolem sequence of Baker, Nowakowski, Shalaby and Sharary \cite{BNSS}).
Section~7 of \cite{ADH} closes by asking for necessary and sufficient conditions for the
existence of $m$-fold Langford sequences with $m \ge 2$. The only
two-fold Langford family printed there (Construction~2.7, defect $6k-1$, order $4k-1$)
lies on the extreme line of the counting bound.

\begin{thmA}
Let $d, l \ge 1$. A two-fold Langford sequence of order $l$ and defect $d$ exists if and
only if $3l \ge 2d-1$.
\end{thmA}

\begin{thmB}
Let $m, d, l \ge 1$.
\begin{enumerate}
\item If an $m$-fold Langford sequence of order $l$ and defect $d$ exists, then
  $(2m-1)l \ge 2d-1$, and if $m$ is odd, $l(2d+l+1) \equiv 0 \pmod 4$.
\item If $(2m-1)l = 2d-1$, an $m$-fold Langford sequence of order $l$ and defect $d$
  exists.
\item An $m$-fold Langford sequence of order $1$ and defect $d$ exists if and only if
  $d \mid m$.
\item No three-fold Langford sequence of order $3$ and defect $6$ exists.
\item For every $m \ge 3$ the conditions of \textup{(i)} are not sufficient: some
  $(d,l)$ satisfies them and has no $m$-fold Langford sequence.
\item If an $m$-fold Langford sequence of order $l$ and defect $d$ exists, then for every
  integer $T \ge 1$, with $r = ml \bmod T$,
  \[
    r(T-r) \;\le\; m \sum_{p=d}^{d+l-1} |p - T| .
  \]
\item If an $m$-fold Langford sequence of order $l$ and defect $d$ exists, then
  $6mld \le 4d^2 + 2m^2l^2 + ml(l-1)$.
\item For every $m \ge 3$ the cell $(d,l) = (3m-3,3)$ satisfies the conditions of
  \textup{(i)} and has no $m$-fold Langford sequence.
\item If $s$ is an $m$-fold Langford sequence of order $l$ and defect $d$, then
  $(2m-1)l = 2d-1$ if and only if $ml < i + s_i$ for every $i \in [1,ml]$, that is, if
  and only if every pair has its left end in $[1,ml]$ and its right end in $[ml+1,2ml]$.
\end{enumerate}
\end{thmB}

At $m = 2$ the parity clause of (i) is vacuous and Theorem~A says that (i) is exact. At
$m = 1$ the conditions of (i) are Simpson's, and exactness there is his theorem
\cite{Simpson}; at $d = 1$ they are the conditions of \cite{BNSS}. Part (v) shows that
nothing of the kind holds for $m \ge 3$, where counting and parity give no complete
existence characterization; explicit witnesses are $(d,l) = (6,3)$ at $m = 3$,
$(m-1,1)$ for even $m \ge 4$ and $(m-2,1)$ for odd $m \ge 5$, and by (viii) they can be
taken in order three for every $m$. Sufficient families do exist for every $m$: the
tight line (ii) and the row $d \mid m$ of (iii). Part (vi) is a family of necessary
conditions, one for every $T$, which contains the counting bound of (i) (take
$T \ge \max(ml, d+l-1)$) and the necessity half of (iii) (take $l = 1$ and $T = d$). Part
(vii) comes from the forced left ends $[1,d]$ and proves (viii) for $m \ge 4$, and (ix)
says that the bound of (i) is attained only when every pair straddles the midpoint.

\subsection*{What is formalized}
Theorem~A and all nine parts of Theorem~B are formalized and machine-checked in Lean~4
with Mathlib \cite{Lean4,Mathlib}, in the accompanying repository \cite{l2lean}, as the
theorems
\begin{center}\small
\begin{tabular}{@{}ll@{}}
\texttt{Langford.twoFold\_exists\_iff} (Theorem A) & \texttt{Langford.order\_one\_iff} (B(iii))\\
\texttt{Langford.necessary} (B(i)) & \texttt{Langford.not\_threeFold\_six\_three} (B(iv))\\
\texttt{Langford.tight\_exists} (B(ii)) & \texttt{Langford.not\_sufficient} (B(v))\\
\texttt{Langford.residue\_bound} (B(vi)) & \texttt{Langford.forced\_endpoint} (B(vii))\\
\texttt{Langford.not\_order\_three} (B(viii)) & \texttt{Langford.tight\_iff\_straddle} (B(ix))
\end{tabular}
\end{center}
about the definition \texttt{Langford.IsLangford} (Definition~\ref{def:seq}). Every numbered
statement below that carries a \emph{Lean} line is proved there under the name given.
Statements that carry a \emph{Status} line instead are proved on paper in this note and
are not formalized, or are finite computations; they are Remark~\ref{rem:nordh},
Corollary~\ref{cor:evenm}, Remark~\ref{rem:residuecert},
Corollary~\ref{cor:residuecases}, Lemma~\ref{lem:I}, Lemma~\ref{lem:Z}, the computations
for $m \ge 3$ in \S\ref{sec:mge3}, and the whole of \S\ref{sec:lp}. The note is the human-readable account
of the formalized proof: it follows the route the formalization takes, so that a reader
can check one against the other without reading Lean.

\subsection*{The shape of the proof of Theorem~A}
Necessity is the distance sum (Theorem~\ref{thm:nec}): the right ends minus the left
ends of the $2l$ pairs sum to $2\sum p$, and no partition of $[1,4l]$ into pairs does
better than $4l^2$. For sufficiency put
\[
  e = 3l-2d+1 \ (\text{the excess}), \qquad
  \delta = d-l, \qquad \mu = 2\delta-1 = l-e .
\]
\begin{itemize}
\item Two sequences $(d,l_1)$ and $(d+l_1,l_2)$ placed one after the other form
  $(d,l_1+l_2)$ (Lemma~\ref{lem:concat}); every in-bound cell with $l \ge 2d$ splits in
  this way into two smaller in-bound cells (Lemma~\ref{lem:split}).
\item Every other in-bound cell with $l \ge 5$ has $-l \le \mu \le l$, and a
  mirror-symmetric ansatz (Lemma~S, Lemma~\ref{lem:S}) turns a permutation $\sigma$ of
  $[1,l]$ whose numbers $\sigma(w)-w+\delta$ and their mirrors tile $[1-l,l]$ --- the
  \emph{signed-permutation problem} $\SP(l,\delta)$ --- into a two-fold sequence of
  defect $l + \delta = d$.
\item Inversion $\sigma \mapsto \sigma^{-1}$ exchanges $\mu$ and $-\mu$; the reversal
  and a two-block reversal $\tau_l$ solve $\mu = 1$ and $\mu = l$; and a composition
  lemma (Lemma~C) carries a solution at $(n,\mu)$ to one at $(n+\mu,\mu)$. So every cell
  descends along its line of fixed $\mu$ to a base cell in the \emph{band}
  $l/2 < \mu \le l$ (Theorem~\ref{thm:R}); one base cell, $(l,\mu) = (4,3)$, is empty, and
  $(7,3)$ is supplied instead.
\item The band is closed, in each residue class of $l$ modulo $4$, by explicit
  block-reversal permutations whose block sizes are affine in $(l,\mu)$: twenty-one
  families and $\tau_l$, each proved by a \emph{class table}, plus eight sporadic cells
  and two small bases (\S\ref{sec:band} and Appendix~\ref{app:tables}).
\item The sixteen in-bound cells with $l \le 4$ are explicit.
\end{itemize}
No part of the argument depends on Simpson's theorem.

\subsection*{Beyond \texorpdfstring{$m = 2$}{m = 2}}
Section~\ref{sec:allm} proves Theorem~B(i)--(iii), (vi), (vii) and (ix) and the
juxtaposition lemma, which gives, with Theorem~A, a sequence at every cell with
$3l \ge 2d-1$ for every even $m$ (Corollary~\ref{cor:evenm}). Section~\ref{sec:mge3}
proves (iv), (v) and (viii), and \S\ref{sec:lp} describes the picture for $m \ge 3$: an
elementary emptiness lemma, a one-parameter family of integer Farkas certificates whose
optimum yields closed-form empty bands with edges
$d/l \to [(2j+1)m - J \mp \sqrt{m^2-J^2}]/(2J)$, $J = j(j+1) < m$, a quadratic
inequality that holds at every integer $T$ without regime hypotheses, an exact finite
test that detects every negative member of the family, realized cells between two empty
bands, and a conjecture. Section~\ref{sec:prior} lists the prior art that was searched.

\section{Definitions, fidelity and certificates}\label{sec:defs}

\begin{definition}\label{def:seq}
Let $m, d, l \ge 1$. A sequence $s = (s_1,\dots,s_{2ml})$ is an \emph{$m$-fold Langford
sequence of order $l$ and defect $d$} if every entry lies in $[d,d+l-1]$ and, for every
$p \in [d,d+l-1]$, there is a set $A_p$ of $m$ positions such that for every $a \in A_p$
\[
  1 \le a, \qquad a + p \le 2ml, \qquad a + p \notin A_p,
\]
and the positions $i \in [1,2ml]$ with $s_i = p$ are exactly the elements of $A_p$ and
of $A_p + p$. For $m = 2$ we say \emph{two-fold}.
\end{definition}
\lean{Langford.IsLangford (Challenge.lean; verbatim copy in L2/Defs.lean)}

The formal definition reads $s$ as a function $\mathbb{N} \to \mathbb{N}$ on the positions
$1,\dots,2ml$ and ignores its other values; intervals are written with bounds. The
hypotheses $d \ge 1$ and $l \ge 1$ in Theorem~A are needed: for $l = 0$ the empty
sequence qualifies, and for $d = 0$ nothing does.

\begin{lemma}[the pair form]\label{lem:pairform}
An $m$-fold Langford sequence of order $l$ and defect $d$ exists if and only if $[1,2ml]$
can be partitioned into $ml$ pairs $\{a,b\}$, $a<b$, in which every $p \in [d,d+l-1]$ is
the difference $b-a$ of exactly $m$ pairs.
\end{lemma}

\begin{proof}
Given $s$ and the sets $A_p$, take the pairs $\{a,a+p\}$, $a \in A_p$. For fixed $p$
these are $m$ pairs with $2m$ distinct endpoints: the $a$ are distinct, so are the
$a+p$, and $a + p \notin A_p$ forbids a left end from being another pair's right end. For
distinct $p$ the endpoint sets are disjoint because a position carries one value. The
$l$ values give $2ml$ endpoints in $[1,2ml]$, so the pairs partition it. Conversely, put
$s_i$ equal to the difference of the pair containing $i$ and $A_p$ equal to the set of
left ends of the pairs of difference $p$.
\end{proof}

\begin{remark}[fidelity: chains]\label{rem:chains}
The definition of \cite{ADH} asks, for each $p$, for ``exactly $2$ disjoint pairs'' of
occurrences at distance $p$, and in a second condition that the two occurrences in each
pair be $p$ apart. The second condition applies within the chosen pairs: read for all
four occurrences it could never hold. The first is read as above: the occurrences of $p$
\emph{can be partitioned} into $m$ pairs at distance $p$, so four occurrences at $a$,
$a+p$, $a+2p$, $a+3p$ (a \emph{chain}) are allowed. This is the $m$-fold wording of
Baker, Nowakowski, Shalaby and Sharary \cite{BNSS}, which \cite{ADH} follows; the
thesis \cite{Alk} prints as a double Langford example the sequence
$(4,2,3,2,4,3,4,2,3,2,4,3)$, whose four $3$s are the chain $3,6,9,12$; and the published
count $1$ of two-fold Skolem sequences of order $1$ \cite{GRS} is the chain $(1,1,1,1)$.
Under the strict reading, which forbids chains, the in-bound cells $(d,l) = (1,1)$,
$(5,4)$ and $(6,6)$ would have no sequence, since each of their sequences contains a
chain, and Theorem~A would be false there. Definition~\ref{def:seq} admits chains.
\end{remark}

\subsection*{Coordinates}
A cell $(d,l)$ is \emph{in-bound} if $3l \ge 2d-1$, equivalently $2d \le 3l+1$. Its
excess is $e = 3l-2d+1 \ge 0$. We use
\[
  \delta = d - l, \qquad \mu = 2\delta - 1 = 2(d-l)-1 = l - e .
\]
In-bound means $\mu \le l$, and $l \le 2d-1$ means $\mu \ge -l$. The line $\mu = l$
($e = 0$) is the tight line; $\mu = -l$ ($e = 2l$) is the line $l = 2d-1$.

\subsection*{Certificates}
The formal development never builds a sequence directly. It builds coloured sets of
integer pairs and proves once that they give a sequence.

\begin{definition}\label{def:cert}
An \emph{$m$-colour certificate} for $(m,d,l)$ is a family $C_0,\dots,C_{m-1}$ of finite
sets of integer pairs $(a,b)$ such that
\begin{enumerate}
\item for each colour $c$, the set $\{b-a : (a,b) \in C_c\}$ equals $[d,d+l-1]$;
\item the set of all endpoints of all pairs of all colours equals $[1,2ml]$;
\item $\sum_c |C_c| \le ml$.
\end{enumerate}
For $m = 2$ we write $(A,B)$ for $(C_0,C_1)$ and call it a two-colour certificate.
\end{definition}
\lean{L2.MultiPairing, L2.TwoFold, L2.TwoFold.toMulti (L2/Multi.lean)}

\begin{lemma}[the bridge]\label{lem:bridge}
If an $m$-colour certificate exists for $(m,d,l)$ with $d,l \ge 1$, then an $m$-fold
Langford sequence of order $l$ and defect $d$ exists.
\end{lemma}

\begin{proof}
Each colour has at least $l$ pairs by (i). With (iii), each has exactly $l$, the colours
are pairwise disjoint, and the difference map is injective on each colour, so each
$p \in [d,d+l-1]$ is the difference of exactly $m$ pairs. The $ml$ pairs have at most
$2ml$ endpoints, and by (ii) they have all of $[1,2ml]$; so no endpoint is shared and
the pairs partition $[1,2ml]$ (a pair $(a,a)$ is impossible because $d \ge 1$). Lemma~\ref{lem:pairform} finishes.
\end{proof}
\lean{L2.MultiPairing.card\_colour, L2.MultiPairing.pairwise\_disjoint,
L2.MultiPairing.card\_fiber, L2.MultiPairing.endpoint\_bijOn,
L2.MultiPairing.isLangford (L2/Bridge.lean)}

The certificate format is the counting device of the formalization \cite{gnlean}
lifted to $m$ colours: families prove only set equalities and a cardinality bound, and
the multiplicity bookkeeping is done in this one lemma. The pair machinery of
\cite{gnlean} (endpoints, differences, translations) is reused, ported to the present
toolchain.

\section{Every multiplicity: the bound, the tight line, the row of order one}\label{sec:allm}

\begin{theorem}[necessity]\label{thm:nec}
Let $m, d, l \ge 1$. If an $m$-fold Langford sequence of order $l$ and defect $d$ exists,
then
\[
  2d + l \le 2ml + 1, \quad\text{that is}\quad (2m-1)l \ge 2d-1,
\]
with equality if and only if every pair has its left end in $[1,ml]$ and its right end in
$[ml+1,2ml]$. If $m$ is odd, moreover $l(2d+l+1) \equiv 0 \pmod 4$.
\end{theorem}

\begin{proof}
Take the pair form, put $H = ml$ and let $L$ and $R$ be the sets of left and right ends;
they partition $[1,2H]$ and have $H$ elements each. Summing $b - a$ over the pairs,
\[
  \sum R - \sum L \;=\; m \sum_{p=d}^{d+l-1} p \;=\; \frac{ml(2d+l-1)}{2}.
\]
An $H$-subset of $[1,2H]$ has sum at most $\sum[H+1,2H]$ and at least $\sum[1,H]$, so
$\sum R - \sum L \le H^2 = m^2l^2$, with equality only if $R = [H+1,2H]$. Hence
$2d + l - 1 \le 2ml$, with the stated equality case. For the parity,
$\sum R + \sum L = H(2H+1)$, so
\[
  4 \sum R \;=\; 2ml(2ml+1) + ml(2d+l-1) \;=\; ml\,(4ml + 2d + l + 1)
  \;\equiv\; ml\,(2d+l+1) \pmod 4,
\]
and the left side is divisible by $4$. For odd $m$, $m$ is invertible modulo $4$.
\end{proof}
\lean{Langford.necessary, L2.necessary\_internal, L2.Nec.sum\_le\_top,
L2.Nec.bottom\_le\_sum, L2.Nec.sum\_right\_sub\_left (L2/Necessity.lean); the equality
clause: Langford.tight\_iff\_straddle, L2.tight\_iff\_straddle\_internal (L2/Straddle.lean)}

The equality case is formalized in the sequence form: for an $m$-fold Langford sequence
$s$ of order $l$ and defect $d$, $2d + l = 2ml + 1$ if and only if $ml < i + s_i$ for
every $i \in [1,ml]$. The two forms agree. If every pair straddles the midpoint, each
$i \le ml$ is a left end and its partner $i + s_i$ exceeds $ml$. Conversely, if
$ml < i + s_i$ for every $i \le ml$ and some $i \le ml$ were a right end, its partner
$a = i - s_i \ge 1$ would be a left end with $s_a = s_i$ and $a + s_a = i \le ml$,
contrary to the hypothesis at $a$; so $[1,ml]$ consists of left ends, and since there are $ml$ of them, the right ends are
$[ml+1,2ml]$. The formal proof uses the same sums: at equality $2\sum L = ml(ml+1)$,
which forces $L = [1,ml]$; conversely, if no right end lies in $[1,ml]$, then
$R = [ml+1,2ml]$ and the distance sum gives $2d + l - 1 = 2ml$.

For even $m$ the parity condition is vacuous (if $l$ is odd, $2d + l + 1$ is even). For
odd $m$ it says: $l \equiv 0 \pmod 4$ allows every $d$, $l \equiv 1 \pmod 4$ needs $d$
odd, $l \equiv 3 \pmod 4$ needs $d$ even, and $l \equiv 2 \pmod 4$ is impossible. These
are the classes of the ordinary case $m = 1$, where the two conditions are Simpson's
\cite{Simpson}; at $d = 1$ they are the conditions of \cite{BNSS} for $m$-fold Skolem
sequences, whose exactness there we take from the literature (the primary text of
\cite{BNSS} was not available to us; its statement agrees in five secondary sources,
among them Table~8 of \cite{ADH}).

\begin{remark}\label{rem:nordh}
Nordh \cite{Nordh} studies \emph{perfect Skolem sets}: a multiset of positive integers
is perfect if $[1,2n]$ can be partitioned into pairs whose differences are exactly its
elements. A two-fold Langford sequence of order $l$ and defect $d$ is precisely a
perfect multiset $2 \times [d,d+l-1]$, and an $m$-fold one a perfect multiset
$m \times [d,d+l-1]$. Nordh's general necessary conditions --- an even number of even
elements, and for every $k$ the $k$ largest elements summing to at most $k(2n-k)$ ---
reduce for $m \times [d,d+l-1]$, $n = ml$, exactly to Theorem~\ref{thm:nec}. Indeed
$f(k) = k(2n-k) - \sum_{i \le k} a_i$ (elements in decreasing order) has increments
decreasing by at least $1$ at every step, so $f$ is concave with $f(0) = 0$ and
$f \ge 0$ everywhere if and only if $f(n) \ge 0$, which is $(2m-1)l \ge 2d-1$; and the
number of even elements, $m$ times the number of even $p \in [d,d+l-1]$, is even for
odd $m$ exactly in the classes listed above. So his conditions add nothing to the
counting bound; he notes himself that they are not sufficient for general multisets.
\end{remark}
\status{proved here; not formalized (it is a comparison with the literature).}

\subsection*{The tight line}
Let $d_0 \ge 1$ and $l = 2d_0 - 1$. Table~1 of \cite{ADH} is the ordinary Langford
sequence of defect $d_0$ and order $l$ with the pairs
\begin{equation}\label{eq:table1}
  (d_0 - r,\; 2d_0 + r) \ (0 \le r \le d_0-1), \qquad
  (2d_0 - 1 - r,\; 3d_0 + r) \ (0 \le r \le d_0-2).
\end{equation}
Their left ends fill $[1,l]$ and their right ends $[l+1,2l]$; the differences
$d_0+2r$ and $d_0+1+2r$ fill $[d_0, 3d_0-2] = [d_0, d_0+l-1]$. As a bijection
$i \mapsto (\text{right end}) - l$ of $[1,l]$ it is the two-block reversal $\tau_l$ of
\S\ref{sec:reduction}: it reverses $[1,d_0]$ and $[d_0+1,l]$.

\begin{theorem}[the tight line, every $m$]\label{thm:tight}
Let $m \ge 1$ and let $l = 2d_0-1$ be odd. Put $d = ((2m-1)l+1)/2 = ml - (d_0-1)$. The
$m$ colours
\[
  C_c = \bigl\{\,(a + cl,\; b + (m-1+c)l) : (a,b) \text{ a pair of } \eqref{eq:table1}\,\bigr\},
  \qquad 0 \le c \le m-1,
\]
form an $m$-colour certificate for $(m,d,l)$. Hence an $m$-fold Langford sequence exists
on the whole line $(2m-1)l = 2d-1$.
\end{theorem}

\begin{proof}
Colour $c$ has left ends $[cl+1,(c+1)l]$ and right ends $[ml+cl+1,ml+(c+1)l]$, so the
union over $c$ is $[1,2ml]$. Its differences are those of \eqref{eq:table1} plus
$(m-1)l$, that is $[d_0+(m-1)l,\, d_0+(m-1)l+l-1] = [d,d+l-1]$. Each colour has $l$
pairs. Lemma~\ref{lem:bridge} finishes. When $l$ is even, $(2m-1)l+1$ is odd and the line
has no integer point.
\end{proof}
\lean{Langford.tight\_exists, L2.table1, L2.tightColour, L2.tight\_multi,
L2.tight\_exists\_internal (L2/Tight.lean, L2/Main.lean)}

By the equality case of Theorem~\ref{thm:nec}, a tight sequence is the same thing as a
permutation $\sigma$ of $[1,ml]$, position $i$ pairing with $ml + \sigma(i)$, whose
displacement multiset $\{\sigma(i) - i\}$ is $m \times [-(l-1)/2,(l-1)/2]$. In the
language of perfect sets this is Theorem~15 of \cite{Nordh} (numbering of
arXiv:math/0506155v1), which should be cited for the bijection; the construction above
is the block-diagonal permutation with $m$ copies of $\tau_l$. At $m = 2$ this closes the
half $d \equiv 2 \pmod 6$ of the tight line that Construction~2.7 of \cite{ADH} does not
reach.

\begin{lemma}[juxtaposition]\label{lem:juxt}
If an $m_1$-fold and an $m_2$-fold Langford sequence of the same order $l$ and defect $d$
exist, so does an $(m_1+m_2)$-fold one: place the second after the first, shifted by
$2m_1l$.
\end{lemma}
\lean{L2.juxtapose, L2.MultiPairing.juxtapose (L2/Multi.lean)}

\begin{corollary}\label{cor:evenm}
For every even $m$, an $m$-fold Langford sequence of order $l$ and defect $d$ exists at
every cell with $3l \ge 2d-1$.
\end{corollary}

\begin{proof}
Juxtapose $m/2$ copies of the two-fold sequence of Theorem~A.
\end{proof}
\status{proved here from Theorem~A and Lemma~\ref{lem:juxt}, both formalized; the
corollary itself is not stated in the development. For odd $m \ge 3$ the same argument
with one ordinary Langford sequence gives every cell with $l \ge 2d-1$ that satisfies the
parity condition, by Simpson's theorem \cite{Simpson}, which is cited and not
re-derived.}

\begin{theorem}[the row of order one]\label{thm:orderone}
Let $m, d \ge 1$. An $m$-fold Langford sequence of order $1$ and defect $d$ exists if and
only if $d \mid m$.
\end{theorem}

\begin{proof}
Here every pair has difference $d$ and the positions are $[1,2m]$. A position $i \le d$
is a left end, since $i - d < 1$. A position $i > d$ is a right end exactly when $i-d$ is
a left end, because the only possible partner of a right end $i$ is $i-d$. By induction
on $i$, the blocks $[1,d]$, $[d+1,2d]$, $[2d+1,3d]$, \dots\ consist alternately of left
ends and right ends. Write $2m = 2dq + r$ with $0 \le r < 2d$. If $1 \le r \le d$, the
position $2m$ lies in a block of left ends and its partner $2m + d$ is out of range. If
$d < r < 2d$, the block $[2dq+1,2dq+d]$ consists of left ends and the partner
$2dq + 2d$ of its last position exceeds $2m$. So $r = 0$ and $d \mid m$. Conversely, if
$m = dk$, the cell $(d,1)$ lies on the tight line of multiplicity $d$
($(2d-1)\cdot 1 = 2d-1$), so Theorem~\ref{thm:tight} gives a $d$-fold sequence, and $k$
copies juxtaposed (Lemma~\ref{lem:juxt}) give an $m$-fold one.
\end{proof}
\lean{Langford.order\_one\_iff, L2.dvd\_of\_order\_one, L2.order\_one\_of\_dvd,
L2.order\_one\_iff\_internal (L2/OrderOne.lean)}

\subsection*{A family of necessary conditions}
For an integer $T$ write
\[
  S(T) \;=\; \sum_{p=d}^{d+l-1} |p - T|
\]
for the total distance from $T$ to the differences.

\begin{theorem}[the residue bound]\label{thm:residue}
Let $m, d, l \ge 1$. If an $m$-fold Langford sequence of order $l$ and defect $d$ exists,
then for every integer $T \ge 1$, with $r = ml \bmod T$,
\[
  r(T-r) \;\le\; m\,S(T).
\]
\end{theorem}

\begin{proof}
Put $n = ml$ and take the pair form. For an integer $y$ let
$w(y) = T - \operatorname{dist}(y, 4T\Z)$, the triangle wave of period $4T$ with values
in $[-T,T]$, and put $V(x) = w(2x - 2n - 1)$ for $x \in \Z$. Three properties of $w$ are
used: it is even; it is $1$-Lipschitz, $|w(y) - w(y')| \le |y - y'|$, as a distance to a
set is; and $w(y + 2T) = -w(y)$, because shifting $y$ by $2T$ turns its distance
$\delta \in [0,2T]$ to $4T\Z$ into $2T - \delta$. Hence $V(x + T) = -V(x)$ and
$|V(x) - V(x')| \le 2|x - x'|$.

\emph{The edge inequality.} Let $(a, a+p)$ be a pair of the sequence. Then
$V(a+p) = -V(a+p-T)$, so
\[
  |V(a) + V(a+p)| \;=\; |V(a) - V(a+p-T)| \;\le\; 2|p - T|,
\]
and for either sign $\sigma \in \{1,-1\}$, $\sigma\,(V(a) + V(a+p)) + 2|p - T| \ge 0$.
Every position of $[1,2n]$ is an end of exactly one pair, and every $p \in [d,d+l-1]$ is
the difference of exactly $m$ pairs, so the sum over the $ml$ pairs is
\begin{equation}\label{eq:residuesum}
  \sigma \sum_{x=1}^{2n} V(x) + 2m\,S(T) \;\ge\; 0 .
\end{equation}

\emph{The total.} As $x$ runs over $[1,2n]$, $2x - 2n - 1$ runs over
$\pm 1, \pm 3, \dots, \pm(2n-1)$, and $w$ is even, so $\sum_{x=1}^{2n} V(x) = 2g(n)$ with
$g(N) = \sum_{k=0}^{N-1} w(2k+1)$. For $0 \le k \le T-1$ one has $0 < 2k+1 < 2T$, so
$w(2k+1) = T - 2k - 1$, and $g(r) = \sum_{k<r} (T - 2k - 1) = r(T-r)$ for
$0 \le r \le T$; in particular $g(T) = 0$, the terms $T-1, T-3, \dots, 1-T$ cancelling.
Since $2(k+T) + 1 = (2k+1) + 2T$, the terms of index $k + T$ are the negatives of those
of index $k$, and $g(N + T) = g(T) - g(N) = -g(N)$. Writing $n = qT + r$ with
$0 \le r < T$,
\[
  \sum_{x=1}^{2n} V(x) \;=\; 2(-1)^q\, r(T-r).
\]

\emph{The sign.} With $\sigma = -(-1)^q$, \eqref{eq:residuesum} reads
$-2r(T-r) + 2mS(T) \ge 0$.
\end{proof}
\lean{Langford.residue\_bound, L2.residue\_bound\_internal (L2/Residue.lean)}

In the formal statement $|p - T|$ is written $(p - T) + (T - p)$ with truncated
subtraction, and $p = d + i$ with $0 \le i < l$. The development follows the proof above:
the triangle wave and its antiperiod, evenness and Lipschitz bound, the block sums
$g$, the total over the positions, and the edge inequality summed over the pairs.

\begin{remark}\label{rem:residuecert}
The potential of the proof is the certificate of Lemma~\ref{lem:familyP}
(\S\ref{sec:lp}): $V(x) = 2W_T(x - c)$ with $c = ml + \tfrac12$, and $(y,z) =
(\sigma V, 2|p - T|)$ is the member with $s = \sigma$. Theorem~\ref{thm:residue} says
that the total of this certificate, with the sign chosen to minimize it, is nonnegative
whenever a sequence exists; it holds for every $T \ge 1$, and in this form it is
formalized. The summation uses only that every position is covered once and every $p$
exactly $m$ times, so, as in Lemma~\ref{lem:farkas}, the bound also holds for fractional
pairings: a cell that violates it has an infeasible linear relaxation.
\end{remark}
\status{proved here; the statement for fractional pairings is not formalized.}

\begin{corollary}[the window form]\label{cor:window}
If an $m$-fold Langford sequence of order $l$ and defect $d$ exists and
$d \le T \le d+l-1$, then with $r = ml \bmod T$, $u = T - d$ and $v = d + l - 1 - T$,
\[
  2r(T-r) \;\le\; m\bigl(u(u+1) + v(v+1)\bigr).
\]
\end{corollary}

\begin{proof}
For $T$ in the window the distances $|p - T|$ are $u, u-1, \dots, 0$ for $p \le T$ and
$1, 2, \dots, v$ for $p > T$, so $2S(T) = u(u+1) + v(v+1)$.
\end{proof}
\lean{L2.residue\_bound\_window (L2/Residue.lean)}

\begin{corollary}[special cases]\label{cor:residuecases}
In Theorem~\ref{thm:residue}:
\begin{enumerate}
\item for $T \ge \max(ml, d+l-1)$ the bound is the counting bound $2d + l \le 2ml + 1$
  of Theorem~\ref{thm:nec};
\item for $l = 1$ and $T = d$ it says $d \mid m$, the necessity half of
  Theorem~\ref{thm:orderone};
\item for $T = d$ it says $r(d-r) \le ml(l-1)/2$, where $r = ml \bmod d$.
\end{enumerate}
\end{corollary}

\begin{proof}
(i) Every $p$ is at most $T$, so $mS(T) = mlT - m\sum_p p$; and $r(T-r) = ml(T - ml)$,
since $r = ml$ if $T > ml$ and both sides vanish if $T = ml$. The bound becomes
$m\sum_p p \le (ml)^2$, that is $ml(2d+l-1)/2 \le (ml)^2$, that is $2d + l - 1 \le 2ml$.
(ii) Here $S(d) = 0$ and $r = m \bmod d$, and $r(d-r) \le 0$ with $0 \le r < d$ forces
$r = 0$. (iii) $S(d) = 0 + 1 + \dots + (l-1) = l(l-1)/2$.
\end{proof}
\status{proved here from Theorem~\ref{thm:residue}; the corollary is not stated in the
development, where the counting bound and the order-one theorem are proved directly
(Theorems~\ref{thm:nec} and~\ref{thm:orderone}).}

\begin{theorem}[the forced-endpoint bound]\label{thm:forced}
Let $m, d, l \ge 1$ and put $e = (2m-1)l - 2d + 1$. If an $m$-fold Langford sequence of
order $l$ and defect $d$ exists, then
\[
  ml\,e \;\le\; (l - 1 + e)^2, \qquad\text{equivalently}\qquad
  6mld \;\le\; 4d^2 + 2m^2l^2 + ml(l-1).
\]
\end{theorem}

\begin{proof}
Put $n = ml$ and let $L$ and $R$ be the sets of left and right ends, as in the proof of
Theorem~\ref{thm:nec}. From $\sum R - \sum L = n(2d+l-1)/2$ and
$\sum R + \sum L = n(2n+1)$,
\[
  4\sum L \;=\; n(4n + 3 - 2d - l).
\]
By Theorem~\ref{thm:nec} and $l \ge 1$, $d \le n$. Every position $x \in [1,d]$ is a left
end, since a right end $a + p$ has $a \ge 1$ and $p \ge d$; and every left end $a$
satisfies $a \le 2n - d$, since $a + p \le 2n$ and $p \ge d$. So $L$ consists of $[1,d]$
and a set $E$ of $n - d$ integers in $[d+1, 2n-d]$, and $\sum E$ is at most the sum of
the $n - d$ largest integers of that interval, $[n+1, 2n-d]$:
\[
  2\sum L \;\le\; d(d+1) + (n-d)(3n-d+1).
\]
With the value of $4\sum L$ this becomes $0 \le 4d^2 - 6nd + 2n^2 + n(l-1)$, the second
form. Since $l - 1 + e = 2(n-d)$, the right side equals $(l - 1 + e)^2 - ne$.
\end{proof}
\lean{Langford.forced\_endpoint, L2.forced\_endpoint\_internal (L2/Forced.lean)}

The development states the second form as $6mld + ml \le 4d^2 + 2(ml)^2 + ml^2$. The
bound is the case $T = d$ of Theorem~\ref{thm:residue}, weakened. Write
$n = ml = qd + r$ with $0 \le r < d$; then
\[
  (n - d)(n - 2d) + r(d-r) \;=\; (1-q)\,d\,\bigl((2-q)d - 2r\bigr) \;\ge\; 0,
\]
since for $q = 0$ both factors are positive, for $q = 1$ the first vanishes, and for
$q \ge 2$ both are at most $0$. With Corollary~\ref{cor:residuecases}(iii),
$ml(l-1)/2 \ge r(d-r) \ge -(n-d)(n-2d)$, and
$2(n-d)(n-2d) + ml(l-1) = 4d^2 - 6nd + 2n^2 + n(l-1)$. In the notation of
\S\ref{sec:lp} the forced-endpoint bound is $2F_1(d) \ge 0$, and the inequality used is
the case $j = 1$, $T = d$ of \eqref{eq:Kbound}. The direct proof above is the one
formalized; it needs only the endpoint sets and their sums.

\section{The reduction at \texorpdfstring{$m = 2$}{m = 2}}\label{sec:reduction}

\subsection*{Concatenation and the unsplittable region}

\begin{lemma}[concatenation]\label{lem:concat}
If two-fold sequences of orders $l_1, l_2$ and defects $d$ and $d + l_1$ exist, then one
of order $l_1 + l_2$ and defect $d$ exists: the second, shifted by $4l_1$, follows the
first.
\end{lemma}

\begin{proof}
The differences are $[d,d+l_1-1] \cup [d+l_1,d+l_1+l_2-1]$, twice each, and the
positions are $[1,4l_1] \cup [4l_1+1,4l_1+4l_2]$. On certificates the two colours are
concatenated colour by colour.
\end{proof}
\lean{L2.TwoFold.concat (L2/Multi.lean)}

\begin{lemma}[splitting]\label{lem:split}
Let $(d,l)$ be in-bound. A split $l = l_1 + l_2$ with both $(d,l_1)$ and $(d+l_1,l_2)$
in-bound exists if and only if $e \ge 5l_1^*$, where $l_1^* = \lceil (2d-1)/3 \rceil =
\lfloor (2d+1)/3 \rfloor$, and then $l_1 = l_1^*$ works. Every in-bound cell with
$l \ge 2d$ splits with $l_1 = l_1^*$, and every cell that does not split has
$l < (16d+2)/9 \le 2d$.
\end{lemma}

\begin{proof}
$(d,l_1)$ is in-bound iff $3l_1 \ge 2d-1$, and $(d+l_1,l-l_1)$ is in-bound iff
$2(d+l_1) \le 3(l-l_1)+1$, that is $5l_1 \le e$. The first is a lower bound on $l_1$ and
the second an upper bound, so a split exists iff the least admissible $l_1^*$ satisfies
the second. If $l \ge 2d$ then $e \ge 4d+1 \ge (10d+5)/3 \ge 5l_1^*$ and
$1 \le l_1^* < l$. If $e < 5l_1^* \le 5(2d+1)/3$ then $9l < 16d+2$.
\end{proof}
\lean{used inside L2.twoFold\_all (L2/Main.lean)}

The unsplittable region $0 \le e < 5\lceil(2d-1)/3\rceil$ is therefore a wedge of slope
roughly $[2/3, 16/9]$ in $l/d$, and the counting bound shows that no concatenation of
consecutive blocks can reach it. The cells with $l \ge 2d-1$ all split except $(d,l) =
(1,1)$ and $(4,7)$: at $l = 2d-1$ one needs $5l_1^* \le 4d-2$, which holds for $d \ge 6$
since $5(2d+1)/3 \le 4d-2$, and for $d \le 5$ fails only at $d = 1$ and $d = 4$. Both
exceptional cells lie on the line $\mu = -l$ and are reached through Lemma~S by
$\tau_l^{-1}$, so Simpson's theorem is not needed anywhere.

\subsection*{The signed-permutation problem}

\begin{definition}\label{def:SP}
Let $l \ge 1$ and $\delta \in \Z$. A permutation $\sigma$ of $[1,l]$ \emph{solves}
$\SP(l,\delta)$ if the $2l$ numbers
\[
  y_w = \sigma(w) - w + \delta \qquad\text{and}\qquad 1 - y_w \qquad (w \in [1,l])
\]
are exactly the elements of $[1-l,l]$. In the formalization $\sigma$ is its graph
$G \subset \Z^2$: the first coordinates of $G$ are $[1,l]$, the second coordinates are
$[1,l]$, $|G| \le l$, and the images of $G$ under the \emph{value} $(w,t) \mapsto
t-w+\delta$ and the \emph{mirror} $(w,t) \mapsto w-t+1-\delta$ have union $[1-l,l]$.
\end{definition}
\lean{L2.SP, L2.val, L2.mir (L2/SP.lean, L2/Rows.lean)}

\begin{lemma}[displacement form]\label{lem:P}
Put $\mu = 2\delta - 1$ and $x_w = 2(\sigma(w)-w) + \mu$. Then $\sigma$ solves
$\SP(l,\delta)$ if and only if $\{|x_w| : w \in [1,l]\} = \{1,3,\dots,2l-1\}$.
\end{lemma}

\begin{proof}
$x_w = 2y_w - 1$, and for an integer $y$ one has $y \in \{k,1-k\}$ exactly when
$|2y-1| = 2k-1$. The set $[1-l,l]$ is the disjoint union of the $l$ pairs
$\{k,1-k\}$, $k \in [1,l]$, and $y_w$, $1-y_w$ always lie in the same pair. So the $2l$
numbers are $[1-l,l]$ iff every pair is hit by exactly one $w$, iff the $|x_w|$ run
through $2k-1$, $k \in [1,l]$, once each.
\end{proof}
\lean{implicit in L2.SP (the set form is the one formalized)}

We call $k = (|x_w|+1)/2$ the \emph{class} of $w$. In the form in which the
construction was first found, a solution was a permutation $\pi$ of $[1,l]$ with a label
$S$ or $D$ on each $u$, the value of $u$ being $u + \pi(u) - 1$ (label $S$) or
$\pi(u) - u$ (label $D$), the $l$ values forming $[\delta,\delta+l-1]$; the two forms
correspond by $u = (|x_w|+1)/2$, $\pi(u) = w$, the label being the sign of $x_w$.

\begin{lemma}[Lemma S]\label{lem:S}
Let $\sigma$ solve $\SP(l,\delta)$ and suppose $d = l + \delta \ge 1$. With
$y_w = \sigma(w) - w + \delta$, the pairs
\[
  A_w = \bigl(l+1-w,\; 2l+y_w\bigr), \qquad B_w = \bigl(2l+1-y_w,\; 3l+w\bigr),
  \qquad w \in [1,l],
\]
form a two-colour certificate $(\{A_w\},\{B_w\})$ for $(2,d,l)$. The resulting two-fold
Langford sequence of order $l$ and defect $d$ is symmetric under $i \mapsto 4l+1-i$,
which exchanges $A_w$ and $B_w$, and every pair joins the outer half
$[1,l] \cup [3l+1,4l]$ to the middle half $[l+1,3l]$.
\end{lemma}

\begin{proof}
The difference of $A_w$ is $2l + y_w - l - 1 + w = d + \sigma(w) - 1$, and so is the
difference of $B_w$. As $\sigma(w)$ runs over $[1,l]$, each colour has the difference set
$[d,d+l-1]$. The left ends of the $A_w$ are $[1,l]$ and the right ends of the $B_w$ are
$[3l+1,4l]$. The remaining endpoints are $2l + y_w$ and $2l + (1-y_w)$, which by the
definition of $\SP$ run over $2l + [1-l,l] = [l+1,3l]$. So the endpoints are $[1,4l]$,
and there are $2l$ pairs. The symmetry is the identity
$4l+1-(2l+y_w) = 2l+1-y_w$ together with $4l+1-(l+1-w) = 3l+w$.
\end{proof}
\lean{L2.colourA, L2.colourB, L2.SP.toTwoFold (L2/LemmaS.lean)}

Not every cell has a sequence of this form. The next invariant shows that exactly one
small cell of each sign is missed.

\begin{lemma}[the value-sum invariant]\label{lem:I}
Let $\sigma$ solve $\SP(l,\delta)$, let $K^+$ be the set of classes $k$ with
$y_w = k$ for some $w$ and $K^-$ the set with $y_w = 1-k$. Then
\[
  \sum_{k \in K^-} (2k-1) \;=\; \frac{l(l-\mu)}{2} \;=\; \frac{le}{2}.
\]
In particular $\SP(4,2)$ ($\mu = 3$) and $\SP(4,-1)$ ($\mu = -3$) have no solution.
\end{lemma}

\begin{proof}
$\sum_w y_w = \sum \sigma(w) - \sum w + l\delta = l\delta$. On the other hand
$K^+ \sqcup K^- = [1,l]$ and $\sum_w y_w = \sum_{K^+} k + \sum_{K^-}(1-k) = l(l+1)/2
- \sum_{K^-}(2k-1)$. So $\sum_{K^-}(2k-1) = l(l+1)/2 - l\delta = l(l-\mu)/2$. At
$(l,\mu) = (4,3)$ this is $2$; at $(4,-3)$ it is $14$, so
$\sum_{K^+}(2k-1) = 16 - 14 = 2$. No set of distinct odd numbers sums to $2$.
\end{proof}
\status{proved here; not formalized. The formal development never needs these two
cells: it uses literal sequences at $(d,l) = (6,4)$ and $(3,4)$ and the base $(7,3)$ in
Theorem~\ref{thm:R}.}

The two cells are $(d,l) = (6,4)$ and $(3,4)$. The first has no mirror-symmetric two-fold
sequence at all (the same invariant, applied after folding a symmetric sequence at its
centre, where no pair can be its own mirror image); it does have the sequence
\[
  (6,9,9,6,8,8,6,7,7,6,9,9,8,8,7,7).
\]
The second has a mirror-symmetric sequence that is
not of the form of Lemma~S. Both are among the sixteen literal cells of
\S\ref{sec:theorem}.

\subsection*{Symmetry, the reversal, \texorpdfstring{$\tau_l$}{tau} and Lemma C}

\begin{lemma}[inversion]\label{lem:inv}
$\sigma$ solves $\SP(l,\delta)$ if and only if $\sigma^{-1}$ solves $\SP(l,1-\delta)$;
in displacement terms $\mu \mapsto -\mu$.
\end{lemma}

\begin{proof}
The graph of $\sigma^{-1}$ consists of the pairs $(t,w)$. The value of $(t,w)$ at
$1-\delta$ is $w-t+1-\delta$, the mirror of $(w,t)$ at $\delta$, and its mirror at
$1-\delta$ is $t-w+\delta$, the value of $(w,t)$ at $\delta$. So the union of values and
mirrors is unchanged.
\end{proof}
\lean{L2.SP.inv, L2.mem\_image\_val\_swap, L2.mem\_image\_mir\_swap}

A second proof uses the reflection $w \mapsto l+1-w$ on the targets with the labels
$S$ and $D$ exchanged; it is a different bijection between the two solution sets.

\begin{lemma}[the reversal and $\tau_l$]\label{lem:rev}
The reversal $w \mapsto l+1-w$ solves $\SP(l,1)$ ($\mu = 1$) for every $l \ge 1$. For
odd $l$, the permutation $\tau_l$ that reverses $[1,(l+1)/2]$ and $[(l+3)/2,l]$ has
displacements $\tau_l(i) - i$ running through $[-(l-1)/2,(l-1)/2]$ once each and solves
$\SP(l,(l+1)/2)$ ($\mu = l$); hence $\tau_l^{-1}$ solves $\mu = -l$.
\end{lemma}

\begin{proof}
For the reversal, $y_w = l+2-2w$ runs through $l, l-2, \dots, 2-l$ and $1-y_w$ through
the integers of the other parity in $[1-l,l-1]$. For $\tau_l$ with $h = (l+1)/2$: on
$[1,h]$ the displacement is $h+1-2i$, running over $h-1, h-3, \dots, 1-h$; on
$[h+1,l]$ it is $l+h+1-2i$, running over $h-2, h-4, \dots, 2-h$. Together they are
$[1-h,h-1]$ once each, so $y = (\text{displacement}) + h$ runs through $[1,l]$ and the
mirrors through $[1-l,0]$.
\end{proof}
\lean{L2.SP.reversal, L2.tauGraph, L2.SP.tau}

\begin{lemma}[Lemma C, with inversion]\label{lem:C}
Let $\sigma_0$ solve $\SP(n,\delta)$ with $\delta \ge 1$, and put $\mu = 2\delta-1$.
Define a permutation $\sigma_1$ of $[1,n+\mu]$ by
\[
  \sigma_1(w) = \sigma_0(w) + \mu \ \ (1 \le w \le n), \qquad
  \sigma_1(n+i) = \tau_\mu(i) \ \ (1 \le i \le \mu).
\]
Then $\sigma_1$ solves $\SP(n+\mu, 1-\delta)$, and $\sigma_1^{-1}$ solves
$\SP(n+\mu,\delta)$.
\end{lemma}

\begin{proof}
The images are $[\mu+1,n+\mu]$ and $[1,\mu]$, so $\sigma_1$ is a permutation. For
$w \le n$, the value of $(w,\sigma_0(w)+\mu)$ at $1-\delta$ is
$\sigma_0(w) - w + \mu + 1 - \delta = \sigma_0(w) - w + \delta$, the old value, and its
mirror is likewise the old mirror; these contribute $[1-n,n]$. For the appended block,
the value of $(n+i,\tau_\mu(i))$ at $1-\delta$ is $(\tau_\mu(i)-i) - n + 1 - \delta$;
the displacements of $\tau_\mu$ are $[1-\delta,\delta-1]$ (Lemma~\ref{lem:rev}), so these
values are $[1-(n+\mu),-n]$ and their mirrors $[n+1,n+\mu]$. The union is
$[1-(n+\mu),n+\mu]$. Lemma~\ref{lem:inv} gives the last claim. No bound on $\mu$ in terms
of $n$ is used.
\end{proof}
\lean{L2.cinvGraph, L2.SP.cinv\_aux, L2.SP.cinv}

\begin{corollary}[descent]\label{cor:descent}
If $\SP(n,\delta)$ has a solution and $\delta \ge 1$, then so has $\SP(n + k\mu,\delta)$
for every $k \ge 0$, where $\mu = 2\delta-1$.
\end{corollary}
\lean{L2.SP.descend}

In its first form Lemma~C carried $(l,\mu)$ to $(l+|\mu|,-\mu)$; composed with inversion
it holds $\mu$ fixed and walks $l$ up the whole line of fixed $\mu$. An equivalent
composite, written with the reflection, gives the same extension.

\subsection*{Reduction to the band}

\begin{theorem}[Theorem R]\label{thm:R}
Assume the band theorem (Theorem~\ref{thm:band}): for every $l \ge 5$ and every odd
$\mu$ with $l < 2\mu \le 2l$, $\SP(l,(\mu+1)/2)$ has a solution. Then $\SP(l,\delta)$ has
a solution for every $l \ge 5$ and every $\delta$ with $-l \le 2\delta-1 \le l$.
\end{theorem}

\begin{proof}
By Lemma~\ref{lem:inv} we may assume $\mu = 2\delta-1 \ge 1$. We use strong induction on
$l$. If $\mu = 1$, use the reversal. If $l < 2\mu$, use the band theorem (it includes
$\mu = l$, which is $\tau_l$). Otherwise $2\mu \le l$; put $l' = l - \mu \ge \mu$ and
obtain a solution at $(l',\mu)$ as follows, then apply Lemma~\ref{lem:C}.
\begin{itemize}
\item If $l' = \mu$: $\tau_\mu$.
\item If $\mu < l' < 2\mu$ and $l' \ge 5$: the band theorem.
\item If $\mu < l' < 2\mu$ and $l' \le 4$: then $\mu$ is odd, $\mu \ge 3$ and
  $\mu < l' \le 4$, so $(l',\mu) = (4,3)$ and $l = 7$. That cell is empty
  (Lemma~\ref{lem:I}), and $(7,3)$ is supplied directly by
  $\sigma = (6,4,7,3,5,2,1)$, whose values $|x_w|$ are $13,7,11,1,3,5,9$.
\item If $l' \ge 2\mu$: then $l' \ge 6$, and the induction hypothesis applies.\qedhere
\end{itemize}
\end{proof}
\lean{L2.cone\_sp, L2.sp\_7\_3 (L2/Cone.lean, L2/Literals.lean)}

Equivalently, every cell descends along its line of fixed $\mu$ to the base
$l_0 = \mu + ((l-\mu) \bmod \mu)$, which lies in $[\mu,2\mu)$, except that the base
$(4,3)$ is replaced by $(7,3)$. The inversion and descent steps are proved in general;
everything specific sits in the band.

\section{The band}\label{sec:band}

The band is the set of cells $(l,\mu)$ with $l \ge 5$, $\mu$ odd and
$l/2 < \mu \le l$. All solutions used there are \emph{block reversals}.

\subsection*{Shapes and the block lemma}
A \emph{shape} is a list of block sizes $n_0,\dots,n_{b-1} \ge 1$ with sum $l$ and a
\emph{target order}, a permutation of the block indices listing them in left-to-right
target position. Block $i$ has source interval $[a_i+1,a_i+n_i]$ with
$a_i = n_0 + \dots + n_{i-1}$ (sources cut left to right) and target interval
$[c_i+1,c_i+n_i]$, where $c_i$ is the sum of the sizes of the blocks placed before $i$ in
the target order; the block maps its source interval onto its target interval reversed:
\[
  \sigma(a_i + r) = c_i + n_i + 1 - r, \qquad 1 \le r \le n_i .
\]
In the formalization block $i$ is the finite set
$\mathrm{row}(a_i,c_i,n_i) = \{(a_i+r,\, c_i+n_i+1-r) : 1 \le r \le n_i\}$ and a family is
a union of rows.

\begin{lemma}[block lemma]\label{lem:block}
Put $M_i = 2(c_i - a_i) + \mu$. On block $i$,
$x_{a_i+r} = M_i + 2n_i + 2 - 4r$ for $1 \le r \le n_i$: an arithmetic progression of
step $4$ from $x_{\mathrm{lo}} = M_i - 2n_i + 2$ to $x_{\mathrm{hi}} = M_i + 2n_i - 2$.
Consequently its classes form a run of one parity and step $2$:
\begin{itemize}
\item type P ($x_{\mathrm{lo}} \ge 1$): classes $(M_i-2n_i+3)/2, \dots, (M_i+2n_i-1)/2$;
\item type N ($x_{\mathrm{hi}} \le -1$): classes $(-M_i-2n_i+3)/2, \dots,
  (-M_i+2n_i-1)/2$;
\item type C ($x_{\mathrm{lo}} \le -1$, $x_{\mathrm{hi}} \ge 1$): if
  $x_{\mathrm{hi}} \equiv 1 \pmod 4$ (type C1), odd classes $1, \dots,
  (x_{\mathrm{hi}}+1)/2$ and even classes $2, \dots, (1-x_{\mathrm{lo}})/2$; if
  $x_{\mathrm{hi}} \equiv 3 \pmod 4$ (type C3), even classes $2, \dots,
  (x_{\mathrm{hi}}+1)/2$ and odd classes $1, \dots, (1-x_{\mathrm{lo}})/2$.
\end{itemize}
\end{lemma}

\begin{proof}
$x_{a_i+r} = 2(c_i+n_i+1-r-a_i-r)+\mu$. The classes follow from $k = (|x|+1)/2$. In the
set form of Definition~\ref{def:SP}, the values of the block run with step $2$ from
$c_i - a_i + 1 - n_i + \delta$ to $c_i - a_i + n_i - 1 + \delta$, and the mirrors from
$a_i - c_i + 2 - n_i - \delta$ to $a_i - c_i + n_i - \delta$.
\end{proof}
\lean{L2.row, L2.mem\_image\_val\_row, L2.mem\_image\_mir\_row, L2.mem\_image\_fst\_row,
L2.mem\_image\_snd\_row (L2/Rows.lean)}

\begin{lemma}[tiling]\label{lem:tiling}
Suppose that the odd class runs of a shape, in increasing order, start at $1$, each
starts $2$ above the end of the previous one, and the last ends at the largest odd number
$\le l$; and that the even runs do the same from $2$ to the largest even number $\le l$.
Then the shape solves $\SP(l,(\mu+1)/2)$.
\end{lemma}

\begin{proof}
The classes $1,\dots,l$ then occur once each, and Lemma~\ref{lem:P} applies; the shape
is a permutation because the blocks tile $[1,l]$ both in source order and in target
order.
\end{proof}

A \emph{family} is a target order together with sizes $n_i$ that are affine functions of
$(l,\mu)$ on a residue class, and a \emph{domain}. Its \emph{class table} lists, for each
block, $n_i$, $M_i$, the type and the class run(s), and the order of the runs in the odd
and in the even chain. A family is proved by four facts, each an identity or a linear
inequality in $(l,\mu)$: (a) the sizes sum to $l$; (b) the chain relations (first start,
abutment, last end) hold identically; (c) the domain is exactly the set where every
$n_i \ge 1$ and every type condition holds; (d) on the residue class every size and run
end is an integer and each C block has a constant residue of $x_{\mathrm{hi}}$ modulo $4$.
Then at every cell of the domain Lemma~\ref{lem:tiling} applies.

The class tables of all families are printed in Appendix~\ref{app:tables}. Each was
derived symbolically from its sizes and target order, and each was re-derived by an
independent implementation: for the even classes of $l$ on every residue sublattice
modulo $32$, for the odd classes with exact linear programming proving the sign
conditions over the domain and the exactness of the domain in both directions. The two
checkers of \S\ref{sec:finite} transcribe every table independently. Every family was,
in addition, built at every domain cell up to at least $l = 400$, and at domain-edge
cells up to $l = 1001$ or beyond, and checked through Lemma~S to a two-fold sequence.

\subsection*{The families}
Table~\ref{tab:families} lists the families. There are twenty-one block-reversal
families and $\tau_l$ (the two-block family TT), twenty-two in all.

\begin{table}[ht]
\centering\small
\caption{The band families. Domains in full, with class tables, in
Appendix~\ref{app:tables}.}\label{tab:families}
\begin{tabular}{@{}lllcl@{}}
\toprule
family & $l$ & $\mu$ & blocks & range of $\mu$ \\
\midrule
FA & $\equiv 2 \ (4)$ & odd & 5 & $(l+4)/2 \le \mu \le (3l-2)/4$ \\
FB2 & $\equiv 2 \ (8)$ & odd & 5 & $(3l+6)/4 \le \mu \le l-1$ \\
FB6 & $\equiv 6 \ (8)$ & odd & 5 & $(3l+10)/4 \le \mu \le l-1$ \\
F4 & $\equiv 6 \ (8)$ & $(3l+2)/4$ & 4 & a line \\
A1 & $\equiv 0 \ (4)$ & $\equiv 1 \ (4)$ & 9 & $l/2+1 \le \mu \le 3l/4-2$ \\
A3 & $\equiv 0 \ (4)$ & $\equiv 3 \ (4)$ & 9 & $l/2+1 \le \mu \le 3l/4-2$, $\mu \le l-5$ \\
U1 & $\equiv 0 \ (4)$ & $\equiv 1 \ (8)$ & 8 & $(2l+11)/3 \le \mu \le l-3$ \\
U3 & $\equiv 0 \ (4)$ & $\equiv 3 \ (8)$ & 9 & $(2l+1)/3 \le \mu \le l-3$ \\
U5 & $\equiv 0 \ (4)$ & $\equiv 5 \ (8)$ & 8 & $(2l+7)/3 \le \mu \le l-3$ \\
U7 & $\equiv 0 \ (4)$ & $\equiv 7 \ (8)$ & 9 & $(4l+11)/5 \le \mu \le l-1$ \\
V & $\equiv 0 \ (4)$ & $\equiv 3 \ (4)$ & 9 & about $(2l+9)/3 \le \mu \le (4l+3)/5$ \\
T & $\equiv 0 \ (4)$ & $\equiv 3 \ (8)$ & 7 & $(8l+11)/9 \le \mu \le l-1$ \\
FA3 & $\equiv 1 \ (4)$ & $\equiv 3 \ (4)$ & 7 & $(l+5)/2 \le \mu \le l-6$ \\
FA1 & $\equiv 1 \ (4)$ & $\equiv 1 \ (4)$ & 8 & $(l+5)/2 \le \mu \le l-4$ \\
FT & $\equiv 1 \ (4)$ & $l-2$ & 6 & a line, $l \ge 9$ \\
FL1 & $\equiv 1 \ (8)$ & $(l+1)/2$ & 7 & a line, $l \ge 17$ \\
FL5 & $\equiv 5 \ (8)$ & $(l+1)/2$ & 6 & a line, $l \ge 13$ \\
R3O & $\equiv 3 \ (4)$ & $\equiv 1 \ (4)$ & 7 & $(l+3)/2 \le \mu \le l-6$ \\
R3E & $\equiv 3 \ (4)$ & $\equiv 3 \ (4)$ & 8 & $(l+7)/2 \le \mu \le l-4$ \\
R3Q & $\equiv 3 \ (4)$ & $\equiv 3 \ (4)$ & 7 & $(l+3)/2 \le \mu \le (2l-1)/3$ \\
R3T1 & $\equiv 3 \ (4)$ & $l-2$ & 6 & a line, $l \ge 7$ \\
TT & odd & $l$ & 2 & $\tau_l$ \\
\bottomrule
\end{tabular}
\end{table}

The families for $l \equiv 0 \pmod 4$ were found by enumerating the block
decompositions of computed solutions and solving, for each solution, the linear system
that its pattern of abutting runs imposes on the sizes; a single solution then yields its
whole two-parameter family. The other classes were found the same way from shape
enumerations. Two structural features recur. In R3O, R3E, R3Q and all the
$l \equiv 1 \pmod 4$ families a \emph{corner} block maps the last $(l-\mu)/2$ sources
reversed onto the first $(l-\mu)/2$ targets, with $M = -l$, and is the only block of
type N. In R3O, R3E and R3Q a \emph{double block}, two blocks consecutive in source and
in target that share one $M$, covers the bottom run of both parities.

\subsection*{Sporadic cells and bases}
Table~\ref{tab:sporadic} lists the eleven explicit solutions. The eight cells with
$l \equiv 0 \pmod 4$ are exactly the band cells of that class that lie in no family
domain; $(5,3)$ and $(9,5)$ complete the class $l \equiv 1 \pmod 4$; $(7,3)$ is the base
of Theorem~\ref{thm:R}.

\begin{table}[ht]
\centering\small
\caption{The explicit solutions, $\sigma = (\sigma(1),\dots,\sigma(l))$, $\delta = (\mu+1)/2$.}
\label{tab:sporadic}
\begin{tabular}{@{}ccl@{}}
\toprule
$(l,\mu)$ & $\delta$ & $\sigma$ \\
\midrule
$(7,3)$ & 2 & $6,4,7,3,5,2,1$ \\
$(5,3)$ & 2 & $3,5,2,4,1$ \\
$(9,5)$ & 3 & $6,5,9,8,1,7,4,3,2$ \\
$(8,5)$ & 3 & $2,6,8,7,5,4,3,1$ \\
$(8,7)$ & 4 & $1,5,4,8,7,2,6,3$ \\
$(12,9)$ & 5 & $7,6,5,1,12,11,10,4,9,8,3,2$ \\
$(12,11)$ & 6 & $2,1,3,9,8,12,11,10,7,6,5,4$ \\
$(16,11)$ & 6 & $1,4,12,11,10,16,15,14,13,9,8,7,6,5,3,2$ \\
$(20,15)$ & 8 & $2,1,14,13,12,11,10,20,19,18,17,16,15,7,6,5,4,3,9,8$ \\
$(24,17)$ & 9 & $2,1,17,16,15,14,13,12,24,23,22,21,20,19,18,9,8,7,6,5,4,3,11,10$ \\
$(32,23)$ & 12 & $2,1,23,22,21,20,19,18,17,16,15,3,32,31,30,29,28,27,26,25,24,$\\
 & & $11,10,9,8,7,6,5,4,14,13,12$ \\
\bottomrule
\end{tabular}
\end{table}
\lean{L2.sp\_7\_3, L2.sp\_5\_3, L2.sp\_9\_5, L2.sp\_8\_5, L2.sp\_8\_7, L2.sp\_12\_9,
L2.sp\_12\_11, L2.sp\_16\_11, L2.sp\_20\_15, L2.sp\_24\_17, L2.sp\_32\_23
(L2/Literals.lean, each by kernel evaluation of a literal finite-set equality)}

\subsection*{The coverage theorems}

\begin{theorem}[$l \equiv 2 \pmod 4$]\label{thm:cov2}
Every band cell with $l \equiv 2 \pmod 4$, $l \ge 6$, has a solution.
\end{theorem}

\begin{proof}
Here $\mu = l$ is even and not in the band, so $\mu$ runs over the odd numbers in
$(l/2,l)$. If $l = 8k+2$ ($k \ge 1$), these are $[4k+3,8k+1]$: FA covers
$[4k+3,6k+1]$ and FB2 covers $[6k+3,8k+1]$. If $l = 8k+6$, they are $[4k+5,8k+5]$: FA
covers $[4k+5,6k+3]$, F4 is $\mu = 6k+5$, and FB6 covers $[6k+7,8k+5]$. In each case the
stated interval is the family's domain restricted to the class.
\end{proof}
\lean{L2.band\_r2 (L2/Band.lean), with L2.famFA\_sp, L2.famFB2\_sp, L2.famFB6\_sp,
L2.famF4\_sp (L2/FamR2.lean)}

\begin{theorem}[$l \equiv 1 \pmod 4$]\label{thm:cov1}
Every band cell with $l \equiv 1 \pmod 4$, $l \ge 5$, has a solution.
\end{theorem}

\begin{proof}
The odd $\mu$ in $(l/2,l]$ are $(l+1)/2$, which is odd, and then $(l+5)/2, \dots, l$,
since $(l+3)/2$ is even. The case $\mu = (l+1)/2$ is FL1 if $l \equiv 1 \pmod 8$,
$l \ge 17$; FL5 if $l \equiv 5 \pmod 8$, $l \ge 13$; and the bases $(5,3)$ and $(9,5)$
at $l = 5$ and $9$. For $(l+5)/2 \le \mu \le l-4$: FA1 if $\mu \equiv 1 \pmod 4$; if
$\mu \equiv 3 \pmod 4$, then $\mu \le l-6$ because $l - 4 \equiv 1 \pmod 4$, and FA3
applies. The case $\mu = l-2$ is FT for $l \ge 9$ and the base $(5,3)$ at $l = 5$. The
case $\mu = l$ is $\tau_l$.
\end{proof}
\lean{L2.band\_r1, with L2.famFA1\_sp, L2.famFA3\_sp, L2.famFT\_sp, L2.famFL1\_sp,
L2.famFL5\_sp (L2/FamR1.lean)}

\begin{theorem}[$l \equiv 3 \pmod 4$]\label{thm:cov3}
Every band cell with $l \equiv 3 \pmod 4$, $l \ge 7$, has a solution.
\end{theorem}

\begin{proof}
Write $l = 4q+3$ and $t = (l-\mu)/2$; the band is $0 \le t \le q$. The case $t = 0$ is
$\tau_l$; $t = 1$ is R3T1; odd $t$ with $3 \le t \le q$ is R3O; even $t$ with
$2 \le t \le q-1$ is R3E. What remains is $t = q$ with $q$ even, so $q \ge 2$; then
$3t = 3q \ge 2q+2$, and R3Q applies (its domain in these coordinates is $t$ even,
$2q+2 \le 3t$, $t \le q$).
\end{proof}
\lean{L2.band\_r3, with L2.famR3O\_sp, L2.famR3E\_sp, L2.famR3Q\_sp, L2.famR3T1\_sp
(L2/FamR3.lean)}

\begin{theorem}[$l \equiv 0 \pmod 4$]\label{thm:cov0}
Every band cell with $l \equiv 0 \pmod 4$, $l \ge 8$, has a solution.
\end{theorem}

\begin{proof}
The band consists of the odd $\mu$ with $l/2 + 1 \le \mu \le l-1$ ($l/2$ is even). For
$l \ge 68$ use the \emph{main intervals}
\[
\begin{array}{llll}
\text{A1, A3:} & [l/2+1,\ 3l/4-2], & \text{U1:} & [(2l+11)/3,\ l-3],\\
\text{U3:} & [(2l+1)/3,\ l-3], & \text{U5:} & [(2l+7)/3,\ l-3],\\
\text{U7:} & [(4l+11)/5,\ l-1], & \text{V:} & [(2l+9)/3,\ (4l+3)/5],\\
\text{T:} & [(8l+11)/9,\ l-1]. & &
\end{array}
\]
Each domain inequality is linear and holds at both ends of the main interval for
$l \ge 68$, so every $\mu$ of the family's class in its main interval lies in its domain.
Split by $\mu \bmod 8$.
\begin{itemize}
\item $\mu \equiv 1$ or $5$: A1, then U1 or U5, up to $l-3$. The intervals overlap
  because $(2l+11)/3 \le 3l/4 - 2$ iff $l \ge 68$ and $(2l+7)/3 \le 3l/4-2$ iff
  $l \ge 52$; and $l - 1$ is never $\equiv 1$ or $5 \pmod 8$.
\item $\mu \equiv 3$: A3, then U3 ($(2l+1)/3 \le 3l/4 - 2$ iff $l \ge 28$) up to $l-3$;
  the cell $\mu = l-1$ (when $l \equiv 4 \pmod 8$) is in T, since
  $(8l+11)/9 \le l-1$ iff $l \ge 20$.
\item $\mu \equiv 7$: A3, then V ($(2l+9)/3 \le 3l/4-2$ iff $l \ge 60$), then U7 up to
  $l-1$. A cell strictly between $(4l+3)/5$ and $(4l+11)/5$ would have
  $5\mu - 4l \in [4,10]$; but $5\mu - 4l \equiv 3 \pmod 8$, since $\mu \equiv 7 \pmod 8$
  and $4l \equiv 0 \pmod{16}$, so there is none.
\end{itemize}
For $8 \le l \le 64$ every band cell was checked against the family domains: all lie in
some domain except $(l,\mu) = (8,5)$, $(8,7)$, $(12,9)$, $(12,11)$, $(16,11)$,
$(20,15)$, $(24,17)$ and $(32,23)$, which are the sporadic cells of
Table~\ref{tab:sporadic}. This finite part is a decision of linear arithmetic for each
of the fifteen values $l/4 \in [2,16]$, and the formal proof makes it exactly so.
\end{proof}
\lean{L2.band\_r0, with L2.famA1\_sp, L2.famA3\_sp, L2.famU1\_sp, L2.famU3\_sp
(L2/FamR0a.lean) and L2.famU5\_sp, L2.famU7\_sp, L2.famV\_sp, L2.famT\_sp
(L2/FamR0b.lean)}

The coverage in this class was also re-proved for every $l$, independently, by exact
lattice geometry (elimination with residue tightening and enumeration of the bounded
remainder), with the same eight sporadic cells as the only uncovered ones.

\begin{theorem}[the band]\label{thm:band}
For every $l \ge 5$ and every odd $\mu$ with $l/2 < \mu \le l$, $\SP(l,(\mu+1)/2)$ has
a solution.
\end{theorem}

\begin{proof}
Split by $l \bmod 4$ and apply Theorems~\ref{thm:cov2}, \ref{thm:cov1},
\ref{thm:cov3} and \ref{thm:cov0}; for even $l$ the band stops at $\mu = l-1$.
\end{proof}
\lean{L2.band\_sp (L2/Band.lean)}

Together with Theorem~\ref{thm:R}, Theorem~\ref{thm:band} gives $\SP(l,\delta)$ on the
whole cone $-l \le 2\delta-1 \le l$ for every $l \ge 5$.

\section{Proof of Theorem A}\label{sec:theorem}

\begin{proof}[Proof of Theorem~A]
Necessity is Theorem~\ref{thm:nec} at $m = 2$: $2d + l \le 4l+1$. For sufficiency we
prove by strong induction on $l$ that every in-bound cell $(d,l)$ has a two-colour
certificate, and conclude by Lemma~\ref{lem:bridge}.
\begin{enumerate}
\item If $l \le 4$, the cell is one of the sixteen in-bound cells of
  Table~\ref{tab:literals}, each an explicit sequence.
\item If $l \ge 5$ and $l \ge 2d$, Lemma~\ref{lem:split} splits it with
  $l_1 = \lfloor(2d+1)/3\rfloor$ into two in-bound cells of smaller order; the induction
  hypothesis gives both, and Lemma~\ref{lem:concat} joins them.
\item If $l \ge 5$ and $l \le 2d-1$, put $\delta = d - l$ and $\mu = 2\delta - 1$. Then
  $\mu \le l$ because the cell is in-bound, and $\mu \ge -l$ because $l \le 2d-1$.
  Theorem~\ref{thm:R} with Theorem~\ref{thm:band} gives a solution of $\SP(l,\delta)$,
  and Lemma~\ref{lem:S} gives a certificate for $(d,l) = (l+\delta,l)$.\qedhere
\end{enumerate}
\end{proof}
\lean{Langford.twoFold\_exists\_iff, L2.twoFold\_all, L2.twoFold\_exists\_iff\_internal
(L2/Main.lean)}

\begin{table}[ht]
\centering\small
\caption{The sixteen in-bound cells with $l \le 4$: one two-fold sequence each, as its
$2l$ pairs of positions.}\label{tab:literals}
\begin{tabular}{@{}cl@{}}
\toprule
$(d,l)$ & pairs \\
\midrule
$(1,1)$ & $[1,2],[3,4]$ \\
$(2,1)$ & $[1,3],[2,4]$ \\
$(1,2)$ & $[1,2],[3,5],[4,6],[7,8]$ \\
$(2,2)$ & $[1,3],[2,5],[4,7],[6,8]$ \\
$(3,2)$ & $[1,4],[2,6],[3,7],[5,8]$ \\
$(1,3)$ & $[1,2],[3,4],[5,7],[6,9],[8,11],[10,12]$ \\
$(2,3)$ & $[1,4],[2,6],[3,7],[5,8],[9,11],[10,12]$ \\
$(3,3)$ & $[1,4],[2,6],[3,8],[5,10],[7,11],[9,12]$ \\
$(4,3)$ & $[1,5],[2,7],[3,9],[4,10],[6,11],[8,12]$ \\
$(5,3)$ & $[1,8],[2,7],[3,9],[4,10],[5,12],[6,11]$ \\
$(1,4)$ & $[1,4],[2,3],[5,7],[6,10],[8,11],[9,13],[12,14],[15,16]$ \\
$(2,4)$ & $[1,3],[2,7],[4,8],[5,9],[6,11],[10,13],[12,15],[14,16]$ \\
$(3,4)$ & $[1,5],[2,6],[3,8],[4,9],[7,13],[10,16],[11,14],[12,15]$ \\
$(4,4)$ & $[1,5],[2,7],[3,9],[4,11],[6,13],[8,14],[10,15],[12,16]$ \\
$(5,4)$ & $[1,6],[2,8],[3,10],[4,12],[5,13],[7,14],[9,15],[11,16]$ \\
$(6,4)$ & $[1,7],[2,11],[3,10],[4,12],[5,13],[6,15],[8,14],[9,16]$ \\
\bottomrule
\end{tabular}
\end{table}

\section{Multiplicity at least three}\label{sec:mge3}

\begin{proposition}\label{prop:sixthree}
There is no three-fold Langford sequence of order $3$ and defect $6$, although
$(m,d,l) = (3,6,3)$ satisfies both conditions of Theorem~\ref{thm:nec}.
\end{proposition}

\begin{proof}
The positions are $[1,18]$ and the differences are $6$, $7$, $8$, three pairs each.
A position $x \le 6$ is a left end, since $x - 6 < 1$; a position $x \ge 13$ is a right
end, since $x + 6 > 18$. As in the proof of Theorem~\ref{thm:nec},
$\sum R - \sum L = 3(6+7+8) = 63$ and $\sum R + \sum L = 171$, so $\sum L = 54$. The
left ends are $[1,6]$, with sum $21$, and three positions of $[7,12]$ with sum $33$;
since $10+11+12 = 33$ is the largest possible such sum, they are $10, 11, 12$. Now $12$
pairs only with $12 + 6 = 18$; $11$ pairs with $17$ or $18$, hence with $17$; $10$ pairs
with $16$, $17$ or $18$, hence with $16$. These are three pairs of difference $6$. But
$7$ is a right end, and its partner $7 - p$ with $p \ge 6$ is a left end in $[1,6]$, so
$7 - p = 1$ and $p = 6$: a fourth pair of difference $6$.
\end{proof}
\lean{Langford.not\_threeFold\_six\_three, L2.not\_threeFold\_six\_three\_internal
(L2/SixThree.lean)}

The conditions hold at $(3,6,3)$: $2 \cdot 6 + 3 = 15 \le 19$ and
$3(12+3+1) = 48 \equiv 0 \pmod 4$.

\begin{theorem}[order three]\label{thm:orderthree}
For every $m \ge 3$ the cell $(d,l) = (3m-3,3)$ satisfies both conditions of
Theorem~\ref{thm:nec}, and no $m$-fold Langford sequence of order $3$ and defect $3m-3$
exists.
\end{theorem}

\begin{proof}
Here $2d + l = 6m - 3 \le 6m + 1 = 2ml + 1$, and for odd $m$,
$l(2d+l+1) = 18m - 6 = 6(3m-1) \equiv 0 \pmod 4$ because $3m - 1$ is even. The excess is
$e = 3(2m-1) - 2(3m-3) + 1 = 4$ and $ml = 3m$, so Theorem~\ref{thm:forced} would give
$12m \le 36$, which fails for $m \ge 4$. At $m = 3$ the cell is $(6,3)$, empty by
Proposition~\ref{prop:sixthree}.
\end{proof}
\lean{Langford.not\_order\_three, L2.not\_order\_three\_internal (L2/Forced.lean)}

The cells $(3m-3,3)$ lie two defects below the tight line of order three, which is
$d = 3m - 1$. For $m = 2$ the cell is $(3,3)$, which has a two-fold sequence (Theorem~A;
Table~\ref{tab:literals}); so among the cells $(3m-3,3)$ with $m \ge 2$ exactly the one
at $m = 2$ is realized.

\begin{proposition}[the top of order two]\label{prop:ordertwotop}
For every $m \ge 3$ there is no $m$-fold Langford sequence of order $2$ and defect
$2m-1$.
\end{proposition}

\begin{proof}
Here $ml = 2m$ and $e = 2(2m-1) - 2(2m-1) + 1 = 1$, so Theorem~\ref{thm:forced} would
give $2m \le 4$. (In the notation of \S\ref{sec:lp}, $F_1(2m-1) = 2 - m < 0$.)
\end{proof}
\lean{L2.not\_order\_two\_top\_internal (L2/Forced.lean)}

The defect $2m-1$ is the largest the counting bound allows in order two
($2d + 2 \le 4m + 1$). For even $m \ge 4$ the cell satisfies both conditions of
Theorem~\ref{thm:nec}, the parity condition being vacuous; for odd $m$ the parity
condition already excludes it ($l \equiv 2 \pmod 4$), so it is a witness only for even
$m$. At $m = 2$ the cell $(3,2)$ is realized (Theorem~A).

\begin{theorem}[counting and parity do not suffice]\label{thm:notsuff}
For every $m \ge 3$ there are $d, l \ge 1$ satisfying the conditions of
Theorem~\ref{thm:nec} with no $m$-fold Langford sequence of order $l$ and defect $d$.
\end{theorem}

\begin{proof}
For $m = 3$ take $(d,l) = (6,3)$ (Proposition~\ref{prop:sixthree}). For even $m \ge 4$
take $(m-1,1)$: it is in-bound, parity is vacuous, and $m-1 \nmid m$. For odd
$m \ge 5$ take $(m-2,1)$: it is in-bound, $1 \cdot (2(m-2)+2) = 2m - 2 \equiv 0
\pmod 4$, and $m - 2 \nmid m$ because $m - 2 \mid m$ would force $m - 2 \mid 2$.
Theorem~\ref{thm:orderone} shows the last two cells empty.
\end{proof}
\lean{Langford.not\_sufficient, L2.not\_sufficient\_internal (L2/OrderOne.lean)}

By Theorem~\ref{thm:orderthree} the witnesses may be taken uniformly in order three,
$(d,l) = (3m-3,3)$ for every $m \ge 3$. Among other small witnesses, $(7,4)$ at $m = 3$
is empty by Lemma~\ref{lem:Z} below and $(3,1)$ at $m = 4$ by
Theorem~\ref{thm:orderone}.

\begin{lemma}[an emptiness lemma]\label{lem:Z}
Let $m \ge 2$. If
\[
  \frac{ml}{2} \;<\; d \;<\; \frac{m(2ml-2l+1)}{3m-2},
\]
then no $m$-fold Langford sequence of order $l$ and defect $d$ exists.
\end{lemma}

\begin{proof}
Put $N = 2ml$, $A = [1,d]$ and $C = [N-d+1,N]$; these are disjoint since $d \le ml$.
Every position of $A$ is a left end and every position of $C$ a right end, and no pair
lies inside $A$ or inside $C$, since every difference is at least $d$. A pair from $A$ to
$C$ has difference at least $N - 2d + 1$, and at most
$t = \max(0, 3d + l - 1 - N)$ values $p \in [d,d+l-1]$ are that large; so at most $mt$
pairs join $A$ to $C$, and at least $2d - 2mt$ positions of $A \cup C$ are paired with
distinct positions of $B = [d+1,N-d]$. Hence $N - 2d \ge 2d - 2mt$, that is
$4d - N \le 2mt$. If $t = 0$ this contradicts $d > ml/2$, which says $4d - N > 0$. If
$t > 0$ it contradicts $d < m(2ml-2l+1)/(3m-2)$, which after expanding says exactly
$4d - N > 2m(3d+l-1-N) = 2mt$.
\end{proof}
\status{proved here; not formalized.}

The interval of Lemma~\ref{lem:Z} has length
$ml(m-2)/(2(3m-2)) + m/(3m-2)$: less than one at $m = 2$, where it contains no integer
(as Theorem~A requires), $3l/14 + 3/7$ at $m = 3$, $2l/5 + 2/5$ at $m = 4$, and about
$ml/6$ only as $m$ grows.

\subsection*{Computations for \texorpdfstring{$m \ge 3$}{m >= 3}}
Every in-bound cell has been decided for $m = 3$ with $l \le 24$, $m = 4$ with $l \le 20$,
and $m = 5, 6$ with $l \le 12$, each by a checked sequence, an exact integer Farkas
certificate (\S\ref{sec:lp}), or the parity condition. In every row there is an empty
band of cells satisfying both conditions of Theorem~\ref{thm:nec}: at $m = 3$ it is
$d \in [11,18]$ at $l = 8$, $[15,28]$ at $l = 12$ and $[30,55]$ at $l = 24$; at $m = 4$,
$[14,27]$ at $l = 8$ and $[34,67]$ at $l = 20$. The last realized $d$ below the band is
about $1.2l$, $1.64l$, $2.1l$ and $2.6l$ for $m = 3,4,5,6$.
\status{checked on the stated ranges; every witness and every certificate is verified
by a standard-library checker with exact arithmetic; not formalized.}

\section{The linear-programming picture for \texorpdfstring{$m \ge 3$}{m >= 3}}\label{sec:lp}

Nothing in this section is formalized; it is proved here, and its proofs and finite
checks were re-derived independently. Fix $(m,d,l)$ and $N = 2ml$.

\begin{lemma}[Farkas certificates]\label{lem:farkas}
Let $y_1,\dots,y_N$ and $z_d,\dots,z_{d+l-1}$ be integers with
$y_a + y_{a+p} + z_p \ge 0$ whenever $1 \le a < a + p \le N$ and $p \in [d,d+l-1]$. If
$\sum_x y_x + m \sum_p z_p < 0$, the cell has no $m$-fold sequence, and its linear
relaxation (fractional pairings) is infeasible.
\end{lemma}

\begin{proof}
Sum the inequality over the $ml$ pairs of a sequence, or over a fractional pairing with
its weights: every position is covered once and every $p$ exactly $m$ times.
\end{proof}

The counting bound is the certificate $y_x = |2x - N - 1|$, $z_p = -2p$, whose total is
$ml\bigl((2m-1)l - 2d + 1\bigr)$.

\begin{lemma}[a family of certificates]\label{lem:familyP}
For an integer $T \ge 1$ and a sign $s = \pm 1$, put $c = (N+1)/2$,
$W_T(u) = T/2 - \operatorname{dist}(u, 2T\Z)$, $g(x) = 2sW_T(x - c)$ and
$z_p = 2|p - T|$. Then $g$ is integer-valued and $y = g$, $z$ satisfy the inequalities of
Lemma~\ref{lem:farkas}. Moreover
\[
  \Phi_s(T) \;=\; \sum_x g(x) + m\sum_p z_p
  \;=\; 2s\,\sigma_2(ml \bmod 2T) + 2m \sum_{p=d}^{d+l-1} |p - T|,
\]
where $\sigma_2(r) = r(T-r)$ for $0 \le r \le T$ and $(r-T)(r-2T)$ for
$T \le r \le 2T$. The counting bound is the member $T \ge \max(ml, d+l-1)$, $s = -1$,
with $\Phi = ml((2m-1)l - 2d + 1)$.
\end{lemma}

\begin{proof}
$2W_T(x-c) = T - \operatorname{dist}(2x-N-1, 4T\Z)$ is an integer. $W_T$ is
$1$-Lipschitz and $W_T(u+T) = -W_T(u)$, so with $u = a - c$,
$g(a) + g(a+p) = 2s[W_T(u) - W_T(u+p-T)] \ge -2|p-T|$. For the closed form, $g$ is
symmetric about $c$, so $\sum g = 4s\sum_{k=1}^{ml} W_T(k-\tfrac12)$; the sum over a
period $2T$ vanishes and the partial sums are $r(T-r)/2$ and $(r-T)(r-2T)/2$.
\end{proof}

Since $|\sigma_2(ml \bmod 2T)| = r(T-r)$ with $r = ml \bmod T$, the smaller of the two
totals is $\min_s \Phi_s(T) = 2H(T)$, where
\[
  H(T) \;=\; mS(T) - r(T-r), \qquad S(T) = \sum_{p=d}^{d+l-1} |p - T| ;
\]
$H(T)$ is half the minimum certificate total, and the statement that it is nonnegative
whenever a sequence exists is the residue bound, Theorem~\ref{thm:residue}, which is
formalized. With Theorem~A it gives the inequality of Proposition~\ref{prop:m2} below,
though not its equality cases.

\begin{theorem}[explicit empty bands]\label{thm:bands}
Let $m \ge 3$, $j \ge 1$ with $J = j(j+1) < m$, and $l \ge 2$. Suppose $d$ is an integer
with
\[
  \bigl|\,2Jd - ((2j+1)m - J)l - J\,\bigr| \;<\;
  \sqrt{(m^2-J^2)\,l^2 - J^2(J+m)/m},
\]
let $T_0$ be an integer nearest to $T^* = m((2j+2)l + 2d - 1)/(2(J+m))$, and suppose
$d \le T_0 \le d+l-1$ and $jT_0 \le ml \le (j+1)T_0$. Then the cell $(m,d,l)$ is empty,
certified by Lemma~\ref{lem:familyP} with $T = T_0$ and $s = (-1)^{j+1}$. For $j = 1$ the
side conditions on $T_0$ follow from $2d \ge (m-2)l + m + 3$.
\end{theorem}

\begin{proof}[Proof outline]
In the regime $jT \le ml \le (j+1)T$, $d \le T \le d+l-1$ one has $\Phi = 2F_j(T)$ with
$F_j(T) = (J+m)T^2 - m((2j+2)l + 2d - 1)T + m^2l^2 + md^2 + m(l-1)d + m(l^2-l)/2$.
Rounding $T^*$ to $T_0$ costs at most $(J+m)/4$. Writing $4(J+m)C - b^2 = mQ_j(d)$ for the
constant $C$ and the linear coefficient $b$ of $F_j$, with
$Q_j(d) = 4Jd^2 - 4((2j+1)ml - Jl + J)d + 4m^2l^2 - (4j+2)ml^2 + 2Jl^2 + (4j+2)ml - 2Jl
- m$, whose discriminant is $16[(m^2-J^2)l^2 + J(J+m)]$, the hypothesis is equivalent
to $mQ_j(d) + (J+m)^2 < 0$, whence $F_j(T_0) < 0$.
\end{proof}
\status{proved here (every algebraic identity also confirmed symbolically); checked
exactly on $3 \le m \le 30$, $2 \le l \le 200$; for $j \ge 2$ the regime conditions are
hypotheses, observed to hold at every tested cell; not formalized.}

\begin{lemma}[the outer bound]\label{lem:outer}
In the regime $jT \le ml \le (j+1)T$, $d\Phi/dT = 4(j+1)(jT - ml) \le 0$ for $T < d$ and
$d\Phi/dT = 4j((j+1)T - ml) \ge 0$ for $T > d + l - 1$. So within a regime the members
with $T \in [d,d+l-1]$, where $\Phi = 2F_j(T)$, dominate the others whenever that window
meets the regime; and if such a member certifies the cell, then $\min_{T \in \mathbb{R}}
F_j(T) < 0$, that is $Q_j(d) < 0$: $d$ lies strictly between the \emph{outer roots}
$[(2j+1)ml - Jl + J \mp \sqrt{(m^2-J^2)l^2 + J(J+m)}]/(2J)$.
\end{lemma}
\status{proved here; not formalized. Every one of the $14{,}975$ cells certified by the
family in the ranges scanned lies strictly inside the outer interval of some $j$ with
positive discriminant (checked).}

Over $m \le 30$, $l \le 200$ the outer radius exceeds the inner radius of
Theorem~\ref{thm:bands} by at most $0.725$. Letting $l \to \infty$:

\begin{corollary}[band edges for large $l$]\label{cor:edges}
For each $j \ge 1$ with $J = j(j+1) < m$, the family of Lemma~\ref{lem:familyP}
certifies an empty band in row $l$ whose edges satisfy
\[
  \frac{d}{l} \;\longrightarrow\; \frac{(2j+1)m - J \mp \sqrt{m^2-J^2}}{2J}
  \qquad (l \to \infty).
\]
\end{corollary}
\status{proved here from Theorem~\ref{thm:bands} and Lemma~\ref{lem:outer}, as a
statement for large $l$; not formalized.}

Table~\ref{tab:edges} gives the limits for small $m$. The first band ends below the
tight line at excess $e/l \to (m - \sqrt{m^2-4})/2$, where $e = (2m-1)l - 2d + 1$; this
is $0.382$, $0.268$, $0.209$, $0.172$ for $m = 3,\dots,6$. The second band ($j = 2$)
exists from $m = 7$ on and the third from $m = 13$ on.

\begin{table}[ht]
\centering\small
\caption{Limits of $d/l$ at the edges of the empty band $j$ (Corollary~\ref{cor:edges}).}
\label{tab:edges}
\begin{tabular}{@{}cccc@{}}
\toprule
$m$ & $j = 1$ & $j = 2$ & $j = 3$ \\
\midrule
3 & $(1.191,\ 2.309)$ & --- & --- \\
4 & $(1.634,\ 3.366)$ & --- & --- \\
5 & $(2.104,\ 4.396)$ & --- & --- \\
6 & $(2.586,\ 5.414)$ & --- & --- \\
7 & $(3.073,\ 6.427)$ & $(2.116,\ 2.717)$ & --- \\
8 & $(3.564,\ 7.436)$ & $(2.392,\ 3.274)$ & --- \\
13 & & $(3.956,\ 5.878)$ & $(3.083,\ 3.500)$ \\
\bottomrule
\end{tabular}
\end{table}

The gap between consecutive bands grows linearly in $l$, roughly as
$(j+1)^2 l/(2m)$ for large $m$. At small $l$ further bands appear: the outer interval of
a $j$ with $J > m$ is nonempty exactly when $l^2 < J/(J-m)$, which accounts for eight
certified cells in the row $l = 1$ outside every $J < m$ interval, such as
$(m,d,l) = (5,2,1)$. At $m = J$ ($m = 2, 6, 12, 20$) a band touches zero at
$d \in \{jl, jl+1\}$, with $\min\Phi = 0$ (checked for $l \le 30$).

On every row decided in \S\ref{sec:mge3} and on three further complete rows, the band
certified by this family equals the set of cells certified by general Farkas
certificates (Lemma~\ref{lem:farkas}): $1{,}219$ cells, none of them realized. All $2{,}840$ emptiness certificates produced in the
computations were re-verified by independent code, and every certified cell with
$N \le 30$ was confirmed empty by exhaustive search.

\begin{proposition}[no obstruction of this kind at $m = 2$]\label{prop:m2}
For $m = 2$, every in-bound $(d,l)$, every integer $T \ge 1$ and both signs,
$\Phi_s(T) \ge 0$. Equality holds only for $s = +1$ with $d = l$, $2T - 3l \in \{-1,0\}$
or $d = l+1$, $2T - 3l \in \{0,1\}$; for $s = -1$ on the tight line with
$T \ge \max(2l, d+l-1)$; and in two degenerate cases ($T = 1$ at $(d,l) = (1,1)$, and
$(d,l,T,s) = (2,1,2,+1)$). So within this family the counting bound is the only member
that ever refutes a cell at $m = 2$.
\end{proposition}

\begin{proof}[Proof outline]
With $ml = 2l$: for $T \ge 2l$, $\sigma_2 = 2l(T-2l) \ge 0$, so $s = +1$ is harmless, and
for $s = -1$, $|p-T| \ge T-p$ gives $\Phi \ge 2l(3l-2d+1) \ge 0$. For $l \le T < 2l$,
$\sigma_2 \le 0$ makes $s = -1$ harmless for $l \ge 2$, and for $s = +1$, $T = l+t$,
$\Phi = 4[\sum_p|p-T| - (l-t)t] \ge 0$ because the distance sum from an integer to $l$
consecutive integers is at least $\lfloor l^2/4 \rfloor$. For $T < l$,
$|\sigma_2| \le T^2/4$ and $\Phi > 0$.
\end{proof}
\status{proved here; checked exactly for $l \le 80$, $T \le 8l$, both signs, by two
implementations; not formalized. It concerns this centred, integer-$T$ family only.}

\begin{lemma}[blow-up]\label{lem:blowup}
If an $m$-fold Langford sequence of order $l_0$ and defect $d_0$ exists, then for every
odd $k$ so does one of order $kl_0$ and defect $kd_0 - (k-1)/2$.
\end{lemma}

\begin{proof}
Replace position $a$ by the block $[k(a-1)+1,ka]$, and each pair $(a,b)$ by the $k$ pairs
$(k(a-1)+i,\ k(b-1)+\tau_k(i))$, $i \in [1,k]$. The displacements of $\tau_k$ are
$[-(k-1)/2,(k-1)/2]$ once each, so a difference $p$ becomes the $k$ differences
$kp + t$, $|t| \le (k-1)/2$, and consecutive $p$ tile
$[kd_0 - (k-1)/2,\ kd_0 - (k-1)/2 + kl_0 - 1]$, each value $m$ times.
\end{proof}
\status{proved here; checked on $264$ instances; not formalized. The same holds for
cells with feasible relaxation. We have not checked the literature for this
multiplication.}

\subsection*{The quadratic inequality without regime hypotheses}
For $j \ge 1$ put
\begin{gather*}
  J = j(j+1), \qquad a = m + J, \qquad b = m\bigl((2j+2)l + 2d - 1\bigr),\\
  C = m^2l^2 + md^2 + m(l-1)d + \frac{ml(l-1)}{2},
\end{gather*}
and $F_j(T) = aT^2 - bT + C$. The coefficients are integers, and $F_j$ is the polynomial
of Theorem~\ref{thm:bands}.

\begin{theorem}[the quadratic inequality]\label{thm:quad}
If an $m$-fold Langford sequence of order $l$ and defect $d$ exists, then
$F_j(T) \ge 0$ for every $j \ge 1$ and every integer $T$.
\end{theorem}

\begin{proof}
Put $n = ml$. For an integer $T \ge 1$ let
\[
  K_j(T) = (n - jT)\bigl(n - (j+1)T\bigr), \qquad
  P(T) = T^2 - (2d+l-1)T + d^2 + (l-1)d + \frac{l(l-1)}{2} ;
\]
expanding, $F_j(T) = K_j(T) + mP(T)$. Two inequalities hold for every integer $T \ge 1$,
whether or not a sequence exists.

First, write $n = qT + r$ with $0 \le r < T$ and put $h = j - q$. Then
\[
  K_j(T) + r(T-r) \;=\; hT\bigl((h+1)T - 2r\bigr),
\]
which is positive for $h \ge 1$, zero for $h = 0$, and for $h \le -1$ a product of two
factors that are at most $0$. So
\begin{equation}\label{eq:Kbound}
  K_j(T) \;\ge\; -r(T-r).
\end{equation}

Second, $P(T) \ge S(T)$. If $d \le T \le d+l-1$ and $u = T - d$, then
$P(T) = u^2 - (l-1)u + l(l-1)/2$, which is $S(T)$ by the closed form in the proof of
Corollary~\ref{cor:window}. If $T < d$, then $P(T) - S(T) = u(u-1)$ with
$u = d - T \ge 1$; if $T > d+l-1$, then $P(T) - S(T) = v(v-1)$ with
$v = T - (d+l-1) \ge 1$. Both are nonnegative because $u$ and $v$ are positive integers.

Hence $F_j(T) \ge mS(T) - r(T-r) = H(T)$ for every integer $T \ge 1$, and $H(T) \ge 0$
by Theorem~\ref{thm:residue}. For $T \le 0$ the three terms $aT^2$, $-bT$ and $C$ are
nonnegative and $C > 0$.
\end{proof}
\status{proved here; not formalized.}

Theorem~\ref{thm:bands} is the special case $T = T_0$ with its regime hypotheses, which
make $\Phi = 2F_j(T_0)$, and with the worst-case rounding cost $a/4$. Theorem~\ref{thm:quad}
needs neither; in particular the regime conditions, which Theorem~\ref{thm:bands} takes
as hypotheses for $j \ge 2$, are not needed for the emptiness of the cell. The
forced-endpoint bound is $2F_1(d) \ge 0$ (Theorem~\ref{thm:forced}). A negative $F_j(T)$
at an integer $T \ge 1$ is a negative $H(T)$, that is, a negative member of the family
of Lemma~\ref{lem:familyP} at the same $T$.

\subsection*{Exact rounding and a finite test}
Put
\[
  t_j = \Bigl\lfloor \frac{b + a}{2a} \Bigr\rfloor, \qquad \rho_j = b - 2at_j,
\]
so that $t_j$ is an integer nearest to $b/(2a)$, ties rounded up, and put
\[
  X = 2Jd - \bigl((2j+1)m - J\bigr)l - J, \qquad D = (m^2 - J^2)l^2 + J(m+J).
\]

\begin{proposition}[exact rounding and a finite cutoff]\label{prop:rounding}
Let $j \ge 1$, with $t_j$, $\rho_j$, $X$ and $D$ as above.
\begin{enumerate}
\item $-a \le \rho_j < a$, $t_j$ minimizes $F_j$ over the integers, and
  $4aF_j(t_j) = 4aC - b^2 + \rho_j^2$. So $F_j(t_j) < 0$ if and only if
  $4aC - b^2 + \rho_j^2 < 0$.
\item $J(4aC - b^2) = m(X^2 - D)$ and $D = (m+J)\bigl(ml^2 - J(l^2-1)\bigr)$. So
  $F_j(t_j) < 0$ if and only if $mX^2 + J\rho_j^2 < mD$.
\item If $F_j(t_j) < 0$, then $J(l^2 - 1) < ml^2$.
\end{enumerate}
\end{proposition}

\begin{proof}
(i) The bounds on $\rho_j$ restate $t_j \le (b+a)/(2a) < t_j + 1$. For every $T$,
$4aF_j(T) = (2aT - b)^2 + 4aC - b^2$, so $F_j$ is smallest over the integers at an integer
nearest to $b/(2a)$, and $2at_j - b = -\rho_j$. (ii) Both identities are checked by
expanding. (iii) If $F_j(t_j) < 0$, then $4aC - b^2 < -\rho_j^2 \le 0$, so by (ii)
$D > X^2 \ge 0$, and $m + J > 0$.
\end{proof}
\status{proved here; checked exactly on the grid described below; not formalized.}

Part (i) is a test in exact integer arithmetic, with no square root.
Theorem~\ref{thm:bands} replaces $\rho_j^2$ by its worst case $a^2$: its hypothesis is
$mX^2 + Ja^2 < mD$, and the outer interval of Lemma~\ref{lem:outer} is $X^2 < D$. For
$l \ge 2$, part (iii) gives $J < ml^2/(l^2-1) \le 4m/3$, so only the indices with
$j(j+1) < 4m/3$, of which there are $O(\sqrt m)$, need to be tested; this counts tests,
each an exact comparison of integers, not bit operations. The cutoff cannot be
simplified to $J < m$, the range of Theorem~\ref{thm:bands}: Example~\ref{ex:19123}
below has $J = 20 > m = 19$. For $l = 1$ the condition in (iii) is vacuous; that row is
settled by Theorem~\ref{thm:orderone}.

\begin{theorem}[completeness for the family]\label{thm:complete}
Let $l \ge 2$. Some member of the family of Lemma~\ref{lem:familyP}, an integer
$T \ge 1$ with a sign $s$, has $\Phi_s(T) < 0$ if and only if the counting bound fails,
$(2m-1)l < 2d - 1$, or $F_j(t_j) < 0$ for some $j \ge 1$ with $J(l^2 - 1) < ml^2$.
\end{theorem}

\begin{proof}
Since $\min_s \Phi_s(T) = 2H(T)$, the left side says that $H(T) < 0$ for some integer
$T \ge 1$. Put $n = ml$ and $e = (2m-1)l - 2d + 1$.

If the counting bound fails, the member $T \ge \max(n, d+l-1)$, $s = -1$, has
$\Phi = ne < 0$. If $F_j(t_j) < 0$, then $t_j \ge 1$, because $F_j > 0$ at the integers
$T \le 0$, and $H(t_j) \le F_j(t_j) < 0$ by the proof of Theorem~\ref{thm:quad}.

Conversely, suppose that the counting bound holds and $H(T) < 0$. By
Proposition~\ref{prop:rounding}(iii) it suffices to find $j \ge 1$ and an integer $T'$
with $F_j(T') < 0$, since then $F_j(t_j) \le F_j(T') < 0$. If $T \ge n$, then
$mS(T) \ge m\sum_p (T - p) = nT - m\sum_p p$ and $r(T-r) = n(T - n)$, so
$H(T) \ge n^2 - m\sum_p p = ne/2 \ge 0$, a contradiction. So $T < n$; put
$j = \lfloor n/T \rfloor \ge 1$. Then $jT \le n < (j+1)T$, so $h = 0$ in the proof of
Theorem~\ref{thm:quad}, and $H(T) = K_j(T) + mS(T)$.
\begin{itemize}
\item If $d \le T \le d+l-1$, then $S(T) = P(T)$ and $F_j(T) = H(T) < 0$.
\item If $T < d$, let $\eta(u) = K_j(u) + m\sum_p (p - u)$ for real $u$. Then
  $\eta(T) = H(T)$, and $\eta'(u) = 2(j+1)(ju - n) \le 0$ for $u \le n/j$. If $n/j < d$,
  then $\eta(n/j) \le \eta(T) < 0$; but $K_j(n/j) = 0$ and every $p$ exceeds $n/j$, so
  $\eta(n/j) > 0$. Hence $T < d \le n/j$, and $F_j(d) = \eta(d) \le \eta(T) < 0$, since
  $\sum_p (p - d) = S(d) = P(d)$.
\item If $T > d+l-1$, let $\eta(u) = K_j(u) + m\sum_p (u - p)$ for real $u$. Then
  $\eta(T) = H(T)$, and $\eta'(u) = 2j((j+1)u - n) \ge 0$ for $u \ge n/(j+1)$. If
  $n/(j+1) > d+l-1$, then $\eta(n/(j+1)) \le \eta(T) < 0$; but $K_j(n/(j+1)) = 0$ and
  every $p$ is below $n/(j+1)$, so $\eta(n/(j+1)) > 0$. Hence
  $n/(j+1) \le d+l-1 < T$, and $F_j(d+l-1) = \eta(d+l-1) \le \eta(T) < 0$.\qedhere
\end{itemize}
\end{proof}
\status{proved here; checked exactly on the grid described below; not formalized.}

The real variable $u$ is used only on intervals where the distance sum is linear in
$u$; the closed form of $S$ inside the window is not used at non-integer points. The
theorem is completeness for this family only: the finite test detects exactly the
negative members of the family of Lemma~\ref{lem:familyP}. It is not completeness for
the linear relaxation, whose infeasibility may have other certificates, and it says
nothing about existence at cells that pass the test; it proves neither
Conjecture~\ref{conj:C} below nor the stronger guess stated after it. For $l = 1$ the
family is exact by Theorem~\ref{thm:orderone}: at $T = d$, $S(d) = 0$ and
$H(d) = -r(d-r)$ with $r = m \bmod d$, which is negative exactly when $d \nmid m$, and
when $d \mid m$ a sequence exists, so by Lemma~\ref{lem:farkas} no member of the family
is negative.

\begin{example}\label{ex:794}
The cell $(m,d,l) = (7,9,4)$ satisfies the counting bound ($22 \le 57$) and the parity
condition ($4 \cdot 23 \equiv 0 \pmod 4$). For $j = 2$,
$(a,b,C,t_2,\rho_2) = (13,287,1582,11,1)$ and $F_2(11) = -2$, so the cell is empty. The
worst-case expression of Theorem~\ref{thm:bands} is $4aC - b^2 + a^2 = 64 > 0$, so its
inner test does not reach this cell, while $4aC - b^2 + \rho_2^2 = -104$. In the form
of Theorem~\ref{thm:residue}, at $T = 11$: $ml \bmod 11 = 6$ and
$7\,(2 + 1 + 0 + 1) = 28 < 6 \cdot (11 - 6) = 30$.
\end{example}

\begin{example}\label{ex:19123}
The cell $(m,d,l) = (19,12,3)$ satisfies the counting bound and the parity condition
($3 \cdot 28 \equiv 0 \pmod 4$). For $j = 4$, so $J = 20 > m$,
$(a,b,C,t_4,\rho_4) = (39,1007,6498,13,-7)$ and $F_4(13) = -2$, so the cell is empty. The
index lies inside the cutoff, $J(l^2-1) = 160 < ml^2 = 171$, and outside the range
$J < m$ of Theorem~\ref{thm:bands}. In the form of Theorem~\ref{thm:residue}, at
$T = 13$: $57 \bmod 13 = 5$ and $19\,(1 + 0 + 1) = 38 < 5 \cdot 8 = 40$.
\end{example}
\status{exact integer arithmetic, repeated by the checker described below; not
formalized.}

The two uniform families of \S\ref{sec:mge3} in this language: at $(d,l) = (3m-3,3)$,
$F_1(3m-2) = -4(m-2)$, negative for every $m \ge 3$, including $m = 3$; at
$(d,l) = (2m-1,2)$, $F_1(2m-1) = 2 - m$. In both, $T$ lies in the window and
$\lfloor ml/T \rfloor = 1$, so these values are $H(T)$.

\emph{Finite checks.} A standard-library checker in exact integer arithmetic tested this
subsection on the grid $1 \le m \le 20$, $1 \le l \le 20$,
$1 \le d \le \lfloor ((2m-1)l+1)/2 \rfloor + 2$, which runs two defects past the counting
line: $42{,}900$ triples. On every triple the finite test (the counting bound, then
$F_j(t_j)$ for the indices within the cutoff) agrees with a direct scan of $H(T)$ over
every integer $T$ with $1 \le T \le \max(ml, d+l-1)$, beyond which $H$ is constant. The
test rejects $28{,}645$ triples that satisfy the counting bound; for the $1{,}026$ of them
with $m \le 10$ and $l \le 8$ the checker rebuilt the integer weights of the certificate,
checked every edge inequality and confirmed the negative total. On the grid, $98$
first-detected obstructions are missed by the worst-case rounding inequality of
Theorem~\ref{thm:bands}; this count includes cells of order one and cells that fail the
parity condition, and is not a count of new empty cells. The checker also verified the
two examples, the identities $4d^2 - 6mld + 2m^2l^2 + ml(l-1) = 2F_1(d) = (l-1+e)^2 - mle$
for $m, l \le 30$ and $d \le ml$, and the identities of the two uniform families for
$3 \le m \le 1000$, and it reproduces the empty bands of \S\ref{sec:mge3} in the rows
$(m,l) = (3,8)$, $(3,12)$, $(3,24)$ and $(4,20)$: $d \in [11,18]$, $[15,28]$, $[30,55]$
and $[34,67]$. A second check decided every cell with $2ml \le 20$ and $m \le 4$, up to
one defect past the counting line, by exhaustive search ($79$ cells, $41$ realized): the
residue bound, for every $T \le 2ml + 2$, and the forced-endpoint bound hold at every
realized cell, the equality case of Theorem~\ref{thm:nec}, in its sequence form, holds
for every sequence found, and among the $38$ empty cells the residue bound rejects $28$ and
the forced-endpoint bound $8$. At $m = 2$, where Theorem~A realizes every in-bound cell,
both bounds hold at every in-bound cell with $l \le 400$, the residue bound for every
$T \le 4l + 2$; the second check also confirmed the two examples.

\subsection*{Realized cells between empty bands}
For $m \ge 7$ a row can contain two empty bands with realized cells between them. In the
row $m = 7$, $l = 12$ the cells $d = 33$ and $d = 38$ are empty by exact certificates
while $d = 34,\dots,37$ are realized. Such islands are witnessed, each by a checked
sequence flanked by certified empty cells in its row, for
$m \in \{7,8,9,11,13,14,15\}$ (for instance $(m,d,l) = (8,34,10)$, $(9,12,3)$,
$(11,20,4)$, $(13,15,4)$, $(14,21,5)$, $(15,18,4)$); and for every even $m \ge 8$ by
Theorem~\ref{thm:orderone} in the row $l = 1$, where $d = m/2$ is realized while
$d = m/2 - 1$ and $d = m/2 + 1$ are not. For other odd $m$ they are not yet witnessed.
For $m \le 6$ the certified band is a single interval in every row checked ($l \le 60$).

\begin{conjecture}[C($m$)]\label{conj:C}
For every $m \ge 3$, an $m$-fold Langford sequence of order $l$ and defect $d$ exists if
and only if the linear relaxation of the cell is feasible and, for odd $m$, the parity
condition holds.
\end{conjecture}

It holds on every decided cell of \S\ref{sec:mge3} and on the islands above. A stronger
guess is that the certificates of Lemma~\ref{lem:familyP} are the only binding ones, so
that the relaxation is infeasible exactly when some $\Phi_s(T) < 0$; they agree on all
$1{,}219$ band cells of the decided rows and, in floating point, on $635$ further cells
with $m \le 14$. With Proposition~\ref{prop:m2}, that guess would make the relaxation
feasible at every in-bound cell for $m = 2$, which Theorem~A implies anyway. No closed
form for the edges of the true empty set is known, and C($m$) is open.

\section{The finite data and the formal development}\label{sec:finite}

\subsection*{The end-to-end certificate}
The certificate folder under \texttt{note/certificates/} of the repository \cite{l2lean}
holds two standard-library Python checkers, \texttt{verify.py} and
\texttt{verify\_second.py}. Each rebuilds, for every $1 \le l \le L$ and every
$1 \le d \le \lfloor(3l+1)/2\rfloor$, a two-fold sequence by the constructions of this
note alone --- concatenation for splittable cells; for the rest, Lemma~S applied to a
permutation obtained by the reversal, $\tau_l$, the descent of Lemma~\ref{lem:C}, or a
band family, with the bases and sporadic cells --- and checks it from the definition
(the positions partition $[1,4l]$; each difference in $[d,d+l-1]$ occurs exactly
twice). The two scripts were written independently, share no code, and transcribe every
family table separately. The first also checks each permutation in displacement form
before lifting, runs every family over its whole domain with $63$ mutated families
rejected, and re-finds the eight sporadic permutations by its own exact search; the
second checks the domain of every family to $l = 160$, compares its closed-form $k$-step
descent with the iterated one on $2{,}701$ chains, and rejects $16$ mutated sequences.
Both also check the $m$-fold tight line for $m \le 6$ and odd $l \le 41$.

The first script ran to $L = 1000$ ($751{,}000$ cells) and the second to $L = 600$, with
no failure; both reproduce at $L = 300$ in fifteen to twenty seconds with
\texttt{python verify.py 300} and \texttt{python verify\_second.py 300}. At $L = 300$
the $67{,}800$ in-bound cells are reached by $25{,}331$ concatenations, $600$ reversals,
$150$ applications of $\tau_l$, $19{,}665$ band-family cells, $22{,}052$ descent chains
and $2$ literal sequences (the cells $(d,l) = (3,4)$ and $(6,4)$, where $\SP$ is empty);
the route counts of the two scripts agree class by class up to the choice between a
literal and a construction at $l \le 4$ and the routing of $(7,3)$.

The finite data establish nothing about all $l$: the certificate is the finite witness of
the theorem, checked to $l = 1000$ by one implementation and to $l = 600$ by a second,
while the statement for all $l$ rests on the class tables, Theorem~R and the induction,
and is the formalized Theorem~A.

\subsection*{Other computations}
Independently of the constructions, every in-bound cell at $m = 2$ with $l \le 54$ was
realized by a constraint solver or by concatenation, and the cells with $l \le 14$ by a
second, exhaustive engine; and the numbers of two-fold
Skolem sequences ($d = 1$) of orders $1,\dots,8$ were recomputed as $1$, $3$, $12$,
$186$, $3212$, $79238$, $2770026$, $127860956$, in agreement with Gr\"uttm\"uller, Rees
and Shalaby \cite{GRS} as tabulated in \cite{Alr}. Before the band families were found,
the signed-permutation problem was solved by exact cover at every cell with
$l \le 100$ except the two empty cells of Lemma~\ref{lem:I}.

\subsection*{The formal development}
The repository \cite{l2lean} pins Lean \texttt{v4.35.0-rc2} and the Mathlib tag of the
same name by a committed manifest. \texttt{Challenge.lean} states the definition and the
ten theorems of \S\ref{sec:intro}; \texttt{Solution.lean} restates them byte for byte and
closes each with the internal theorem of the same name in the library \texttt{L2}, which
has these modules:
\begin{itemize}\setlength{\itemsep}{0pt plus 1pt}
\item \texttt{Defs} --- the verbatim copy of \texttt{Langford.IsLangford};
\item \texttt{Pairs} --- endpoints, differences and translations of finite sets of
  pairs, ported from \cite{gnlean};
\item \texttt{Rows} --- the row of a block reversal, its values and mirrors
  (Lemma~\ref{lem:block});
\item \texttt{SP} --- $\SP$ in graph form, inversion, the reversal, $\tau_l$, Lemma~C
  and the descent (\S\ref{sec:reduction});
\item \texttt{FamR0a}, \texttt{FamR0b}, \texttt{FamR1}, \texttt{FamR2}, \texttt{FamR3}
  --- one definition and one theorem per family, each the class table of
  Appendix~\ref{app:tables} with the audited domain as hypotheses;
\item \texttt{Literals} --- the eleven permutations of Table~\ref{tab:sporadic} and the
  sixteen cells of Table~\ref{tab:literals};
\item \texttt{Band}, \texttt{Cone} --- Theorems~\ref{thm:cov2}--\ref{thm:band} and
  Theorem~\ref{thm:R};
\item \texttt{Multi}, \texttt{LemmaS}, \texttt{Bridge} --- certificates, concatenation,
  juxtaposition, Lemma~S and Lemma~\ref{lem:bridge};
\item \texttt{Tight}, \texttt{Necessity}, \texttt{OrderOne}, \texttt{SixThree} ---
  Theorems~\ref{thm:tight}, \ref{thm:nec}, \ref{thm:orderone}, \ref{thm:notsuff} and
  Proposition~\ref{prop:sixthree};
\item \texttt{Residue} --- the residue bound, Theorem~\ref{thm:residue} and
  Corollary~\ref{cor:window};
\item \texttt{Forced} --- the forced-endpoint bound, the order-three cells and the top
  of order two, Theorems~\ref{thm:forced} and~\ref{thm:orderthree} and
  Proposition~\ref{prop:ordertwotop};
\item \texttt{Straddle} --- the rigidity of the counting bound, the equality clause of
  Theorem~\ref{thm:nec};
\item \texttt{Main} --- the induction of \S\ref{sec:theorem} and the internal theorems.
\end{itemize}
Internals work over $\Z$, where linear arithmetic with negative slopes is decided by
Lean's \texttt{omega}; the pinned statements are over $\mathbb{N}$, and the casts happen
in \texttt{Bridge} and \texttt{Main}, and inside the proof of the residue bound, whose
potential takes negative values. Finite facts --- the eleven permutations, the
sixteen literal cells, and in the proof for $(6,3)$ the three-element subsets of $[7,12]$
with sum $33$ --- are decided by the kernel on literal finite sets; there is no other
search inside the proofs.
A build-time audit (\texttt{Test/Axioms.lean}) walks every declaration of the development
--- \texttt{\detokenize{DESK_FILLS}} constants in about \texttt{\detokenize{DESK_FILLS}}
lines of Lean --- and admits only the axioms
\texttt{propext}, \texttt{Classical.choice} and \texttt{Quot.sound}; a source guard (\texttt{scripts/check-source.py}) rejects
\texttt{sorry}, \texttt{axiom}, \texttt{native\_decide}, \texttt{implemented\_by},
\texttt{extern} and the kernel-bypass options. A standard-library script,
\texttt{scripts/check\_langford.py}, recomputes the finite content independently of
Lean: the sixteen literal cells and the eleven permutations from the definitions, every
family at every domain cell with $l \le 120$ lifted through Lemma~S, the $m$-fold tight
line for $m \le 6$ and odd $l \le 21$, the row $l = 1$ for $m \le 12$, the emptiness of
$(3,6,3)$ by exhaustive search, the bounds of Theorems~\ref{thm:residue}
and~\ref{thm:forced} against brute-force existence on small cells, and two controls that
must fail.

\section{Prior art}\label{sec:prior}

On 2026-09-26 searches were made of the full texts of \cite{ADH} (arXiv version and
journal version) and of the thesis \cite{Alk}; of seven further theses from the same
department; of Nordh's papers on perfect Skolem sets; of arXiv, Crossref, zbMATH Open,
the On-Line Encyclopedia of Integer Sequences, a public repository of formalized
conjectures, and the Palomar registry of machine-checked results; and of the works citing
\cite{ADH}, \cite{BNSS} and Nordh's papers, and the bibliography of the survey
\cite{FM}. No statement of necessary and sufficient conditions for two-fold or $m$-fold
Langford sequences with $d \ge 2$ was located. The question is printed as open in
\cite{ADH} (\S7, identical in both versions) and twice in \cite{Alk}; Table~8 of
\cite{ADH}, which summarizes known existence results, has $m$-fold Skolem rows and no
$m$-fold Langford row. The search was bounded, and on 2026-09-27 its blind spots were
re-examined. The body of the survey \cite{FM} and the Handbook chapter on Skolem-type
sequences could be searched only through full-text indexes, which contain no $m$-fold
Langford statement; the 2026 paper of Zhang and Cao on generalized signed Langford
sequences (its abstract and references) and the 2017 paper of L\'opez and Muntaner-Batle on
distance labelings (read in full) concern different objects --- a signed Langford
sequence is, by Nordh's Theorem~15, a one-fold sequence with a different difference set;
Nowakowski's 1975 master's thesis and Dillon's 1973 paper on the generalized
Langford--Skolem problem remain unread, and MathSciNet was not searched. Failure to
retrieve a paper is not evidence about its contents, and nothing here is a claim of
priority beyond that wording.

The residue bound and the forced-endpoint bound were the subject of a targeted search on
2026-09-27 (the same sources, with the terms weight function, linear programming and
Farkas); no condition of this kind was located, and Nordh's parity and density
conditions (Remark~\ref{rem:nordh}) are the only comparable ones, reducing to the
counting bound. No claim of novelty is made beyond that wording.

Four connections follow.

\emph{Perfect Skolem sets.} As Remark~\ref{rem:nordh} explains, the objects of this
note are Nordh's perfect multisets $m \times [d,d+l-1]$ \cite{Nordh}; his Theorem~15 is
the tight-line bijection, and his necessary conditions reduce to Theorem~\ref{thm:nec}.
Nordh's introduction notes that the existence question for general multisets is NP-complete; the multisets here are highly
structured, and that result is context only.

\emph{Constructions in \cite{ADH}.} Construction~2.7 there is a two-fold Langford family
on the tight line, and the sequences obtained by concatenating two Langford sequences,
or interlacing a hooked one with its reverse, exist exactly when $l \ge 2d-1$ or
$(d,l) \in \{(2,1),(2,2)\}$. Nothing in print constructs two-fold Langford sequences
inside the unsplittable region.

\emph{Chains and multiplicity.} The generalized Langford sequences of multiplicity $s$
(the copies of each value in one chain with equal gaps; see \cite{MMNS}) are a different
object. A chain of four copies contains two disjoint pairs, so such sequences with $s = 4$
are two-fold Langford sequences, but their span forces $l \ge 3d-2$, far from the
unsplittable region.

\emph{Defect one.} The $m$-fold Skolem sequences of \cite{BNSS}, and their descendants
(near, hooked, extended and disjoint variants), are all of defect $1$. The exact counts
of two-fold Skolem sequences of \cite{GRS} are reproduced by the search (\S\ref{sec:finite}).

In the registry of machine-checked results (read to 2026-09-27), Langford sequences appear
only as sources of the formalization \cite{gnlean} and of its successor; the repository of
formalized conjectures contains no Langford or Skolem statement.

\section*{Provenance}

This note and the accompanying Lean development were produced with Claude (Anthropic),
running with parallel Opus subagents; the mathematics,
the searches and the checkers arose in that process. JD Jones owns and directed the
project and is responsible for the submission. No independent human review or
endorsement by the cited authors is claimed.

\appendix

\section{The class tables}\label{app:tables}

Conventions (\S\ref{sec:band}). Sizes, $M_i = 2(c_i - a_i) + \mu$ and the class runs are
affine in $(l,\mu)$ on the family's residue class; for the line families they are written
in $l$ alone. The target order lists the block indices from left to right in target
position. The type is P, N, C1 or C3 as in Lemma~\ref{lem:block}; a C block contributes
two runs, $b_i^{+}$ (its positive $x$) and $b_i^{-}$ (its negative $x$). The two chains
list the runs of each parity in increasing order; in each chain the first run starts at
$1$ or $2$, each run starts $2$ above the end of the previous one, and the last ends at
$l$ or $l-1$, all identically in $(l,\mu)$. Every domain below is the set of cells of
the residue class with $\mu$ odd at which every size is at least $1$ and every type
condition holds. The tables were generated from the sizes and orders of the audited
transcription and every identity was checked symbolically; every family was also built
at every domain cell with $l \le 240$ and checked in displacement form.

% family FA
\famtable{FA}{$l\equiv 2 \pmod 4$}{2,3,0,4,1}{$(l+4)/2 \le \mu \le (3l-2)/4$}{\texttt{famFA}}{
$b_{0}$ & $(\mu-1)/2$ & $2l-\mu$ & P & odd $[l-\mu+2,\ l-1]$ \\
$b_{1}$ & $(l+2)/4$ & $3l/2$ & P & even $[(l+2)/2,\ l]$ \\
$b_{2}$ & $(3l-4\mu+2)/4$ & $-l/2$ & N & odd $[(2\mu-l+2)/2,\ l-\mu]$ \\
$b_{3}$ & $(l-2)/4$ & $-l/2$ & N & even $[2,\ (l-2)/2]$ \\
$b_{4}$ & $(2\mu-l)/4$ & $(2\mu-l-2)/2$ & P & odd $[1,\ (2\mu-l-2)/2]$ \\
}{$b_{4}$, $b_{2}$, $b_{0}$}{$b_{3}$, $b_{1}$}

% family FB2
\famtable{FB2}{$l\equiv 2 \pmod 8$}{1,3,0,4,2}{$(3l+6)/4 \le \mu \le l-1$ (so $l \ge 10$)}{\texttt{famFB2}}{
$b_{0}$ & $(4\mu-l-2)/8$ & $(9l-4\mu-2)/4$ & P & odd $[(5l-4\mu+6)/4,\ l-1]$ \\
$b_{1}$ & $(l-2)/8$ & $(l+2)/4$ & P & even $[2,\ (l-2)/4]$ \\
$b_{2}$ & $(3l+2)/8$ & $(5l+2)/4$ & P & even $[(l+6)/4,\ l]$ \\
$b_{3}$ & $l-\mu$ & $-l/2$ & N & odd $[(4\mu-3l+6)/4,\ (5l-4\mu-2)/4]$ \\
$b_{4}$ & $(4\mu-3l+2)/8$ & $(4\mu-3l-2)/4$ & P & odd $[1,\ (4\mu-3l-2)/4]$ \\
}{$b_{4}$, $b_{3}$, $b_{0}$}{$b_{1}$, $b_{2}$}

% family FB6
\famtable{FB6}{$l\equiv 6 \pmod 8$}{1,3,0,4,2}{$(3l+10)/4 \le \mu \le l-1$}{\texttt{famFB6}}{
$b_{0}$ & $(4\mu-l+2)/8$ & $(9l-4\mu+2)/4$ & P & even $[(5l-4\mu+6)/4,\ l]$ \\
$b_{1}$ & $(l+2)/8$ & $(l-2)/4$ & P & odd $[1,\ (l-2)/4]$ \\
$b_{2}$ & $(3l-2)/8$ & $(5l-2)/4$ & P & odd $[(l+6)/4,\ l-1]$ \\
$b_{3}$ & $l-\mu$ & $-l/2$ & N & even $[(4\mu-3l+6)/4,\ (5l-4\mu-2)/4]$ \\
$b_{4}$ & $(4\mu-3l-2)/8$ & $(4\mu-3l+2)/4$ & P & even $[2,\ (4\mu-3l-2)/4]$ \\
}{$b_{1}$, $b_{2}$}{$b_{4}$, $b_{3}$, $b_{0}$}

% family F4
\famtable{F4}{$l\equiv 6 \pmod 8$, $\mu = (3l+2)/4$}{2,0,3,1}{$l \ge 6$}{\texttt{famF4}}{
$b_{0}$ & $(3l-2)/8$ & $(5l-2)/4$ & P & odd $[(l+6)/4,\ l-1]$ \\
$b_{1}$ & $(l+2)/4$ & $3l/2$ & P & even $[(l+2)/2,\ l]$ \\
$b_{2}$ & $(l-2)/4$ & $-l/2$ & N & even $[2,\ (l-2)/2]$ \\
$b_{3}$ & $(l+2)/8$ & $(l-2)/4$ & P & odd $[1,\ (l-2)/4]$ \\
}{$b_{3}$, $b_{0}$}{$b_{2}$, $b_{1}$}

% family A1
\famtable{A1}{$l\equiv 0 \pmod 4$, $\mu \equiv 1 \pmod 4$}{8,3,6,7,2,0,5,1,4}{$l \ge 8$, $\mu \ge 5$, $l/2+1 \le \mu \le 3l/4-2$}{\texttt{famA1}}{
$b_{0}$ & $(\mu-1)/4$ & $(4l-\mu-1)/2$ & P & odd $[(2l-\mu+3)/2,\ l-1]$ \\
$b_{1}$ & $1$ & $2l-1$ & P & even $\{l\}$ \\
$b_{2}$ & $(l-\mu+1)/4$ & $(3l-\mu-5)/2$ & P & even $[l/2,\ (2l-\mu-3)/2]$ \\
$b_{3}$ & $(2\mu-l+2)/4$ & $(2\mu-l)/2$ & P & odd $[1,\ (2\mu-l)/2]$ \\
$b_{4}$ & $(\mu-1)/4$ & $(4l-\mu-5)/2$ & P & even $[(2l-\mu+1)/2,\ l-2]$ \\
$b_{5}$ & $(\mu-1)/4$ & $(4l-3\mu-7)/2$ & P & odd $[l-\mu,\ (2l-\mu-5)/2]$ \\
$b_{6}$ & $(3l-4\mu-4)/4$ & $(2-l)/2$ & N & odd $[(2\mu-l+4)/2,\ l-\mu-2]$ \\
$b_{7}$ & $(l-4)/4$ & $(2-l)/2$ & N & even $[2,\ (l-4)/2]$ \\
$b_{8}$ & $1$ & $\mu-2l+2$ & N & odd $\{(2l-\mu-1)/2\}$ \\
}{$b_{3}$, $b_{6}$, $b_{5}$, $b_{8}$, $b_{0}$}{$b_{7}$, $b_{2}$, $b_{4}$, $b_{1}$}

% family A3
\famtable{A3}{$l\equiv 0 \pmod 4$, $\mu \equiv 3 \pmod 4$}{8,3,6,7,2,0,5,1,4}{$l \ge 8$, $\mu \ge 7$, $l/2+1 \le \mu \le 3l/4-2$, $\mu \le l-5$}{\texttt{famA3}}{
$b_{0}$ & $(\mu+1)/4$ & $(4l-\mu-3)/2$ & P & odd $[(2l-\mu+1)/2,\ l-1]$ \\
$b_{1}$ & $1$ & $2l-1$ & P & even $\{l\}$ \\
$b_{2}$ & $(l-\mu-1)/4$ & $(3l-\mu-7)/2$ & P & even $[l/2,\ (2l-\mu-5)/2]$ \\
$b_{3}$ & $(2\mu-l+2)/4$ & $(2\mu-l)/2$ & P & odd $[1,\ (2\mu-l)/2]$ \\
$b_{4}$ & $(\mu-3)/4$ & $(4l-\mu-3)/2$ & P & even $[(2l-\mu+3)/2,\ l-2]$ \\
$b_{5}$ & $(\mu+1)/4$ & $(4l-3\mu-5)/2$ & P & odd $[l-\mu,\ (2l-\mu-3)/2]$ \\
$b_{6}$ & $(3l-4\mu-4)/4$ & $(2-l)/2$ & N & odd $[(2\mu-l+4)/2,\ l-\mu-2]$ \\
$b_{7}$ & $(l-4)/4$ & $(2-l)/2$ & N & even $[2,\ (l-4)/2]$ \\
$b_{8}$ & $1$ & $\mu-2l+2$ & N & even $\{(2l-\mu-1)/2\}$ \\
}{$b_{3}$, $b_{6}$, $b_{5}$, $b_{0}$}{$b_{7}$, $b_{2}$, $b_{8}$, $b_{4}$, $b_{1}$}

% family U1
\famtable{U1}{$l\equiv 0 \pmod 4$, $\mu \equiv 1 \pmod 8$}{5,7,1,3,4,0,2,6}{$3\mu \ge 2l+11$, $5 \le \mu \le l-3$, $\mu \le 2l-15$}{\texttt{famU1}}{
$b_{0}$ & $(2l-\mu+1)/8$ & $(6l+\mu-5)/4$ & P & odd $[(2l+\mu+3)/4,\ l-1]$ \\
$b_{1}$ & $(3\mu-2l-3)/8$ & $(2l+\mu+3)/4$ & P & even $[(2l-\mu+9)/4,\ (\mu-1)/2]$ \\
$b_{2}$ & $(\mu+3)/4$ & $(4l-\mu-1)/2$ & P & even $[(2l-\mu+1)/2,\ l]$ \\
$b_{3}$ & $(2l-\mu+1)/8$ & $(2l-\mu-3)/4$ & P & odd $[1,\ (2l-\mu-3)/4]$ \\
$b_{4}$ & $(2l-\mu-7)/8$ & $(2l-\mu-3)/4$ & P & even $[2,\ (2l-\mu-7)/4]$ \\
$b_{5}$ & $1$ & $(\mu-2l+1)/2$ & N & even $\{(2l-\mu+1)/4\}$ \\
$b_{6}$ & $(\mu-1)/4$ & $l-1$ & P & odd $[(2l-\mu+5)/4,\ (2l+\mu-5)/4]$ \\
$b_{7}$ & $(l-\mu-1)/2$ & $1-l$ & N & even $[(\mu+3)/2,\ (2l-\mu-3)/2]$ \\
}{$b_{3}$, $b_{6}$, $b_{0}$}{$b_{4}$, $b_{5}$, $b_{1}$, $b_{7}$, $b_{2}$}

% family U3
\famtable{U3}{$l\equiv 0 \pmod 4$, $\mu \equiv 3 \pmod 8$}{5,8,0,3,6,4,1,7,2}{$3\mu \ge 2l+1$, $7 \le \mu \le l-3$, $\mu \le 2l-29$}{\texttt{famU3}}{
$b_{0}$ & $(\mu-3)/4$ & $l+1$ & P & even $[(2l-\mu+11)/4,\ (2l+\mu-3)/4]$ \\
$b_{1}$ & $(\mu+1)/4$ & $(4l-\mu-3)/2$ & P & odd $[(2l-\mu+1)/2,\ l-1]$ \\
$b_{2}$ & $(2l-\mu+3)/8$ & $(6l+\mu+1)/4$ & P & even $[(2l+\mu+5)/4,\ l]$ \\
$b_{3}$ & $(2l-\mu-21)/8$ & $(2l-\mu-1)/4$ & P & even $[4,\ (2l-\mu-13)/4]$ \\
$b_{4}$ & $1$ & $(2l-\mu+1)/2$ & P & even $\{(2l-\mu+3)/4\}$ \\
$b_{5}$ & $1$ & $(\mu-2l+7)/2$ & N & even $\{(2l-\mu-5)/4\}$ \\
$b_{6}$ & $(2l-\mu+3)/8$ & $(2l-\mu-17)/4$ & C1 & odd $[1,\ (2l-\mu-9)/4]$; even $\{2\}$ \\
$b_{7}$ & $(3\mu-2l+7)/8$ & $(2l+\mu-7)/4$ & P & odd $[(2l-\mu-1)/4,\ (\mu-1)/2]$ \\
$b_{8}$ & $(l-\mu-1)/2$ & $1-l$ & N & odd $[(\mu+3)/2,\ (2l-\mu-3)/2]$ \\
}{$b_{6}^{+}$, $b_{7}$, $b_{8}$, $b_{1}$}{$b_{6}^{-}$, $b_{3}$, $b_{5}$, $b_{4}$, $b_{0}$, $b_{2}$}

% family U5
\famtable{U5}{$l\equiv 0 \pmod 4$, $\mu \equiv 5 \pmod 8$}{4,7,0,3,5,1,6,2}{$3\mu \ge 2l+7$, $5 \le \mu \le l-3$, $\mu \le 2l-11$}{\texttt{famU5}}{
$b_{0}$ & $(\mu-1)/4$ & $l+1$ & P & odd $[(2l-\mu+9)/4,\ (2l+\mu-1)/4]$ \\
$b_{1}$ & $(\mu+3)/4$ & $(4l-\mu-1)/2$ & P & even $[(2l-\mu+1)/2,\ l]$ \\
$b_{2}$ & $(2l-\mu-3)/8$ & $(6l+\mu-1)/4$ & P & odd $[(2l+\mu+7)/4,\ l-1]$ \\
$b_{3}$ & $(2l-\mu-3)/8$ & $(2l-\mu+1)/4$ & P & even $[2,\ (2l-\mu-3)/4]$ \\
$b_{4}$ & $1$ & $(\mu-2l+1)/2$ & N & odd $\{(2l-\mu+1)/4\}$ \\
$b_{5}$ & $(2l-\mu-3)/8$ & $(2l-\mu-7)/4$ & P & odd $[1,\ (2l-\mu-7)/4]$ \\
$b_{6}$ & $(3\mu-2l+1)/8$ & $(2l+\mu-1)/4$ & P & even $[(2l-\mu+5)/4,\ (\mu-1)/2]$ \\
$b_{7}$ & $(l-\mu-1)/2$ & $1-l$ & N & even $[(\mu+3)/2,\ (2l-\mu-3)/2]$ \\
}{$b_{5}$, $b_{4}$, $b_{0}$, $b_{2}$}{$b_{3}$, $b_{6}$, $b_{7}$, $b_{1}$}

% family U7
\famtable{U7}{$l\equiv 0 \pmod 4$, $\mu \equiv 7 \pmod 8$}{1,6,2,4,7,0,8,3,5}{$5\mu \ge 4l+11$, $7 \le \mu \le l-1$, $\mu \le 2l-5$, $3\mu \le 4l+1$}{\texttt{famU7}}{
$b_{0}$ & $(\mu+1)/8$ & $(8l-\mu+3)/4$ & P & even $[(4l-\mu+7)/4,\ l]$ \\
$b_{1}$ & $(\mu+1)/8$ & $(3\mu-1)/4$ & P & odd $[(\mu+5)/4,\ (\mu-1)/2]$ \\
$b_{2}$ & $(\mu+1)/8$ & $(8l-5\mu-1)/4$ & P & even $[(4l-3\mu+5)/4,\ (2l-\mu-1)/2]$ \\
$b_{3}$ & $(2l-\mu-1)/4$ & $(2l+\mu-1)/2$ & P & odd $[(\mu+3)/2,\ l-1]$ \\
$b_{4}$ & $(5\mu-4l-3)/8$ & $(4l-3\mu+1)/4$ & P & odd $[l-\mu+2,\ (\mu-3)/4]$ \\
$b_{5}$ & $(\mu+1)/8$ & $(8l-3\mu+1)/4$ & P & even $[(2l-\mu+3)/2,\ (4l-\mu-1)/4]$ \\
$b_{6}$ & $l-\mu$ & $(1-\mu)/2$ & N & even $[(5\mu-4l+5)/4,\ (4l-3\mu-3)/4]$ \\
$b_{7}$ & $(l-\mu+1)/2$ & $\mu-l$ & N & odd $[1,\ l-\mu]$ \\
$b_{8}$ & $(5\mu-4l-3)/8$ & $(5\mu-4l+1)/4$ & P & even $[2,\ (5\mu-4l-3)/4]$ \\
}{$b_{7}$, $b_{4}$, $b_{1}$, $b_{3}$}{$b_{8}$, $b_{6}$, $b_{2}$, $b_{5}$, $b_{0}$}

% family V
\famtable{V}{$l\equiv 0 \pmod 4$, $\mu \equiv 3 \pmod 4$}{8,0,5,7,3,4,1,2,6}{$5\mu \ge 3l+11$, $3\mu \ge 2l+9$, $2\mu \ge l+6$, $7 \le \mu \le l-5$, $5\mu \le 4l+3$}{\texttt{famV}}{
$b_{0}$ & $(l-\mu+3)/4$ & $(l+\mu-1)/2$ & P & even $[(\mu+1)/2,\ l/2]$ \\
$b_{1}$ & $(\mu-3)/4$ & $(4l-\mu+1)/2$ & P & odd $[(2l-\mu+5)/2,\ l-1]$ \\
$b_{2}$ & $(\mu+1)/4$ & $(4l-\mu+1)/2$ & P & even $[(2l-\mu+3)/2,\ l]$ \\
$b_{3}$ & $(4l-5\mu+7)/4$ & $(2l-\mu+1)/2$ & P & odd $[(2\mu-l)/2,\ (3l-3\mu+3)/2]$ \\
$b_{4}$ & $(3\mu-2l-5)/4$ & $(2l-\mu+1)/2$ & P & even $[l-\mu+3,\ (\mu-3)/2]$ \\
$b_{5}$ & $(5\mu-3l-7)/4$ & $(\mu-l-1)/2$ & C3 & even $[2,\ (3\mu-2l-5)/2]$; odd $[1,\ (2\mu-l-4)/2]$ \\
$b_{6}$ & $(2\mu-l-2)/4$ & $(5l-4\mu+6)/2$ & P & odd $[(3l-3\mu+7)/2,\ (2l-\mu+1)/2]$ \\
$b_{7}$ & $(4l-5\mu+7)/4$ & $(1-\mu)/2$ & N & even $[(3\mu-2l-1)/2,\ l-\mu+1]$ \\
$b_{8}$ & $(l-\mu-1)/4$ & $(\mu-3l-1)/2$ & N & even $[(l+4)/2,\ (2l-\mu-1)/2]$ \\
}{$b_{5}^{-}$, $b_{3}$, $b_{6}$, $b_{1}$}{$b_{5}^{+}$, $b_{7}$, $b_{4}$, $b_{0}$, $b_{8}$, $b_{2}$}

% family T
\famtable{T}{$l\equiv 0 \pmod 4$, $\mu \equiv 3 \pmod 8$}{2,3,6,0,4,1,5}{$9\mu \ge 8l+11$, $11 \le \mu \le l-1$, $3\mu \le 4l-3$, $7\mu \le 8l-3$}{\texttt{famT}}{
$b_{0}$ & $(\mu+1)/4$ & $(4l-\mu+1)/2$ & P & even $[(2l-\mu+3)/2,\ l]$ \\
$b_{1}$ & $(\mu-3)/8$ & $(8l-\mu-1)/4$ & P & odd $[(4l-\mu+7)/4,\ l-1]$ \\
$b_{2}$ & $(8l-7\mu+5)/8$ & $(\mu+1)/4$ & C3 & even $[2,\ (4l-3\mu+1)/4]$; odd $[1,\ l-\mu]$ \\
$b_{3}$ & $(9\mu-8l-3)/8$ & $(\mu+1)/4$ & P & odd $[l-\mu+2,\ (5\mu-4l-3)/4]$ \\
$b_{4}$ & $(\mu-3)/8$ & $(8l-5\mu+3)/4$ & P & even $[(4l-3\mu+9)/4,\ (2l-\mu-1)/2]$ \\
$b_{5}$ & $(\mu+1)/4$ & $2l-\mu$ & P & odd $[(4l-3\mu+5)/4,\ (4l-\mu-1)/4]$ \\
$b_{6}$ & $l-\mu$ & $(1-\mu)/2$ & N & odd $[(5\mu-4l+5)/4,\ (4l-3\mu-3)/4]$ \\
}{$b_{2}^{-}$, $b_{3}$, $b_{6}$, $b_{5}$, $b_{1}$}{$b_{2}^{+}$, $b_{4}$, $b_{0}$}

% family FA3
\famtable{FA3}{$l\equiv 1 \pmod 4$, $\mu \equiv 3 \pmod 4$}{6,1,3,4,0,5,2}{$(l+5)/2 \le \mu \le l-6$}{\texttt{famFA3}}{
$b_{0}$ & $(\mu+1)/4$ & $(4l-\mu+1)/2$ & P & odd $[(2l-\mu+3)/2,\ l]$ \\
$b_{1}$ & $(\mu+1)/4$ & $(2l-\mu-1)/2$ & P & even $[(l-\mu+2)/2,\ (l-1)/2]$ \\
$b_{2}$ & $(l-1)/4$ & $(3l-1)/2$ & P & even $[(l+3)/2,\ l-1]$ \\
$b_{3}$ & $(l-\mu-2)/4$ & $(l-\mu)/2$ & P & even $[2,\ (l-\mu-2)/2]$ \\
$b_{4}$ & $(l-\mu+2)/4$ & $(l-\mu)/2$ & P & odd $[1,\ (l-\mu)/2]$ \\
$b_{5}$ & $(2\mu-l-1)/4$ & $(l+1)/2$ & P & odd $[(l-\mu+4)/2,\ (\mu-1)/2]$ \\
$b_{6}$ & $(l-\mu)/2$ & $-l$ & N & odd $[(\mu+3)/2,\ (2l-\mu-1)/2]$ \\
}{$b_{4}$, $b_{5}$, $b_{6}$, $b_{0}$}{$b_{3}$, $b_{1}$, $b_{2}$}

% family FA1
\famtable{FA1}{$l\equiv 1 \pmod 4$, $\mu \equiv 1 \pmod 4$}{7,1,3,4,0,5,2,6}{$(l+5)/2 \le \mu \le l-4$}{\texttt{famFA1}}{
$b_{0}$ & $(\mu-1)/4$ & $(4l-\mu-1)/2$ & P & even $[(2l-\mu+3)/2,\ l-1]$ \\
$b_{1}$ & $(\mu-5)/4$ & $(2l-\mu+1)/2$ & P & odd $[(l-\mu+6)/2,\ (l-3)/2]$ \\
$b_{2}$ & $(l-1)/4$ & $(3l+3)/2$ & P & odd $[(l+5)/2,\ l]$ \\
$b_{3}$ & $(l-\mu)/4$ & $(l-\mu+2)/2$ & P & even $[2,\ (l-\mu)/2]$ \\
$b_{4}$ & $(l-\mu+4)/4$ & $(l-\mu+2)/2$ & P & odd $[1,\ (l-\mu+2)/2]$ \\
$b_{5}$ & $(2\mu-l-1)/4$ & $(l+1)/2$ & P & even $[(l-\mu+4)/2,\ (\mu-1)/2]$ \\
$b_{6}$ & $1$ & $l$ & P & odd $\{(l+1)/2\}$ \\
$b_{7}$ & $(l-\mu)/2$ & $-l$ & N & even $[(\mu+3)/2,\ (2l-\mu-1)/2]$ \\
}{$b_{4}$, $b_{1}$, $b_{6}$, $b_{2}$}{$b_{3}$, $b_{5}$, $b_{7}$, $b_{0}$}

% family FT
\famtable{FT}{$l\equiv 1 \pmod 4$, $\mu = l-2$}{5,1,3,0,4,2}{$l \ge 9$}{\texttt{famFT}}{
$b_{0}$ & $(l-1)/4$ & $(3l-1)/2$ & P & even $[(l+3)/2,\ l-1]$ \\
$b_{1}$ & $(l-5)/4$ & $(l+1)/2$ & P & odd $[3,\ (l-3)/2]$ \\
$b_{2}$ & $(l-1)/4$ & $(3l+3)/2$ & P & odd $[(l+5)/2,\ l]$ \\
$b_{3}$ & $1$ & $1$ & P & odd $\{1\}$ \\
$b_{4}$ & $(l-1)/4$ & $(l+1)/2$ & P & even $[2,\ (l-1)/2]$ \\
$b_{5}$ & $1$ & $-l$ & N & odd $\{(l+1)/2\}$ \\
}{$b_{3}$, $b_{1}$, $b_{5}$, $b_{2}$}{$b_{4}$, $b_{0}$}

% family FL1
\famtable{FL1}{$l\equiv 1 \pmod 8$, $\mu = (l+1)/2$}{6,0,2,4,5,1,3}{$l \ge 17$}{\texttt{famFL1}}{
$b_{0}$ & $1$ & $l$ & P & odd $\{(l+1)/2\}$ \\
$b_{1}$ & $(l-1)/8$ & $(7l-3)/4$ & P & even $[(3l+5)/4,\ l-1]$ \\
$b_{2}$ & $(l-9)/8$ & $(3l+1)/4$ & P & odd $[(l+11)/4,\ (l-3)/2]$ \\
$b_{3}$ & $(l-1)/4$ & $(3l+3)/2$ & P & odd $[(l+5)/2,\ l]$ \\
$b_{4}$ & $(l-1)/8$ & $(l+3)/4$ & P & even $[2,\ (l-1)/4]$ \\
$b_{5}$ & $(l+7)/8$ & $(l+3)/4$ & P & odd $[1,\ (l+3)/4]$ \\
$b_{6}$ & $(l-1)/4$ & $-l$ & N & even $[(l+7)/4,\ (3l-3)/4]$ \\
}{$b_{5}$, $b_{2}$, $b_{0}$, $b_{3}$}{$b_{4}$, $b_{6}$, $b_{1}$}

% family FL5
\famtable{FL5}{$l\equiv 5 \pmod 8$, $\mu = (l+1)/2$}{5,1,3,4,0,2}{$l \ge 13$}{\texttt{famFL5}}{
$b_{0}$ & $(l+3)/8$ & $(7l+1)/4$ & P & odd $[(3l+5)/4,\ l]$ \\
$b_{1}$ & $(l+3)/8$ & $(3l-3)/4$ & P & even $[(l+3)/4,\ (l-1)/2]$ \\
$b_{2}$ & $(l-1)/4$ & $(3l-1)/2$ & P & even $[(l+3)/2,\ l-1]$ \\
$b_{3}$ & $(l-5)/8$ & $(l-1)/4$ & P & even $[2,\ (l-5)/4]$ \\
$b_{4}$ & $(l+3)/8$ & $(l-1)/4$ & P & odd $[1,\ (l-1)/4]$ \\
$b_{5}$ & $(l-1)/4$ & $-l$ & N & odd $[(l+7)/4,\ (3l-3)/4]$ \\
}{$b_{4}$, $b_{5}$, $b_{0}$}{$b_{3}$, $b_{1}$, $b_{2}$}

% family R3O
\famtable{R3O}{$l\equiv 3 \pmod 4$, $\mu \equiv 1 \pmod 4$}{6,1,3,4,0,5,2}{$(l+3)/2 \le \mu \le l-6$}{\texttt{famR3O}}{
$b_{0}$ & $(\mu-1)/4$ & $(4l-\mu-1)/2$ & P & even $[(2l-\mu+3)/2,\ l-1]$ \\
$b_{1}$ & $(\mu-1)/4$ & $(2l-\mu+1)/2$ & P & odd $[(l-\mu+4)/2,\ (l-1)/2]$ \\
$b_{2}$ & $(l+1)/4$ & $(3l+1)/2$ & P & odd $[(l+3)/2,\ l]$ \\
$b_{3}$ & $(l-\mu-2)/4$ & $(l-\mu)/2$ & P & even $[2,\ (l-\mu-2)/2]$ \\
$b_{4}$ & $(l-\mu+2)/4$ & $(l-\mu)/2$ & P & odd $[1,\ (l-\mu)/2]$ \\
$b_{5}$ & $(2\mu-l+1)/4$ & $(l-1)/2$ & P & even $[(l-\mu+2)/2,\ (\mu-1)/2]$ \\
$b_{6}$ & $(l-\mu)/2$ & $-l$ & N & even $[(\mu+3)/2,\ (2l-\mu-1)/2]$ \\
}{$b_{4}$, $b_{1}$, $b_{2}$}{$b_{3}$, $b_{5}$, $b_{6}$, $b_{0}$}

% family R3E
\famtable{R3E}{$l\equiv 3 \pmod 4$, $\mu \equiv 3 \pmod 4$}{7,0,2,4,5,1,6,3}{$(l+7)/2 \le \mu \le l-4$}{\texttt{famR3E}}{
$b_{0}$ & $1$ & $l$ & P & even $\{(l+1)/2\}$ \\
$b_{1}$ & $(\mu+1)/4$ & $(4l-\mu+1)/2$ & P & odd $[(2l-\mu+3)/2,\ l]$ \\
$b_{2}$ & $(\mu-3)/4$ & $(2l-\mu-1)/2$ & P & even $[(l-\mu+4)/2,\ (l-3)/2]$ \\
$b_{3}$ & $(l-3)/4$ & $(3l+1)/2$ & P & even $[(l+5)/2,\ l-1]$ \\
$b_{4}$ & $(l-\mu)/4$ & $(l-\mu+2)/2$ & P & even $[2,\ (l-\mu)/2]$ \\
$b_{5}$ & $(l-\mu+4)/4$ & $(l-\mu+2)/2$ & P & odd $[1,\ (l-\mu+2)/2]$ \\
$b_{6}$ & $(2\mu-l-3)/4$ & $(l+3)/2$ & P & odd $[(l-\mu+6)/2,\ (\mu-1)/2]$ \\
$b_{7}$ & $(l-\mu)/2$ & $-l$ & N & odd $[(\mu+3)/2,\ (2l-\mu-1)/2]$ \\
}{$b_{5}$, $b_{6}$, $b_{7}$, $b_{1}$}{$b_{4}$, $b_{2}$, $b_{0}$, $b_{3}$}

% family R3Q
\famtable{R3Q}{$l\equiv 3 \pmod 4$, $\mu \equiv 3 \pmod 4$}{6,0,2,4,5,1,3}{$(l+3)/2 \le \mu \le (2l-1)/3$}{\texttt{famR3Q}}{
$b_{0}$ & $(2\mu-l-1)/2$ & $l$ & P & even $[l-\mu+2,\ \mu-1]$ \\
$b_{1}$ & $(\mu+1)/4$ & $(4l-\mu+1)/2$ & P & odd $[(2l-\mu+3)/2,\ l]$ \\
$b_{2}$ & $(2l-3\mu+3)/4$ & $(2l-\mu-1)/2$ & P & even $[(\mu+1)/2,\ l-\mu]$ \\
$b_{3}$ & $(l-\mu)/2$ & $l+\mu-1$ & P & even $[\mu+1,\ l-1]$ \\
$b_{4}$ & $(\mu-3)/4$ & $(\mu-1)/2$ & P & even $[2,\ (\mu-3)/2]$ \\
$b_{5}$ & $(\mu+1)/4$ & $(\mu-1)/2$ & P & odd $[1,\ (\mu-1)/2]$ \\
$b_{6}$ & $(l-\mu)/2$ & $-l$ & N & odd $[(\mu+3)/2,\ (2l-\mu-1)/2]$ \\
}{$b_{5}$, $b_{6}$, $b_{1}$}{$b_{4}$, $b_{2}$, $b_{0}$, $b_{3}$}

% family R3T1
\famtable{R3T1}{$l\equiv 3 \pmod 4$, $\mu = l-2$}{5,1,3,0,4,2}{$l \ge 7$}{\texttt{famR3T1}}{
$b_{0}$ & $(l-3)/4$ & $(3l+1)/2$ & P & even $[(l+5)/2,\ l-1]$ \\
$b_{1}$ & $(l-3)/4$ & $(l+3)/2$ & P & odd $[3,\ (l-1)/2]$ \\
$b_{2}$ & $(l+1)/4$ & $(3l+1)/2$ & P & odd $[(l+3)/2,\ l]$ \\
$b_{3}$ & $1$ & $1$ & P & odd $\{1\}$ \\
$b_{4}$ & $(l-3)/4$ & $(l-1)/2$ & P & even $[2,\ (l-3)/2]$ \\
$b_{5}$ & $1$ & $-l$ & N & even $\{(l+1)/2\}$ \\
}{$b_{3}$, $b_{1}$, $b_{2}$}{$b_{4}$, $b_{5}$, $b_{0}$}

% family TT
\famtable{TT}{$l$ odd, $\mu = l$}{0,1}{$l \ge 3$ (this is $\tau_l$)}{\texttt{tauGraph}}{
$b_{0}$ & $(l+1)/2$ & $l$ & P & odd $[1,\ l]$ \\
$b_{1}$ & $(l-1)/2$ & $l$ & P & even $[2,\ l-1]$ \\
}{$b_{0}$}{$b_{1}$}


\begin{thebibliography}{99}

\bibitem{ADH}
A.~H.~Alkasasbeh, D.~Dyer and J.~Howell,
\emph{Graceful labellings of variable windmills using Skolem-type sequences},
Ars Combin. \textbf{159} (2024), 109--131; doi:10.61091/ars159-11.
Preprint version: \emph{Graceful labellings of variable windmills using Skolem
sequences}, arXiv:2112.04265v1 (2021).

\bibitem{Alk}
A.~H.~Alkasasbeh,
\emph{Graceful labellings of new families of windmill and snake graphs},
PhD thesis, Memorial University of Newfoundland, 2021; handle 20.500.14783/2068.

\bibitem{Alr}
F.~Alruhaymi,
\emph{Mathematical and algorithmic methods for finding disjoint Rosa-type sequences},
MSc thesis, Memorial University of Newfoundland, 2018, Table~4.2.

\bibitem{BNSS}
C.~Baker, R.~Nowakowski, N.~Shalaby and A.~Sharary,
\emph{M-fold and extended m-fold Skolem sequences},
Utilitas Math. \textbf{45} (1994), 153--167; Zbl~0808.05001.

\bibitem{FM}
N.~Franceti\'c and E.~Mendelsohn,
\emph{A survey of Skolem-type sequences and Rosa's use of them},
Math. Slovaca \textbf{59} (2009), 39--76; doi:10.2478/s12175-008-0110-3.

\bibitem{GRS}
M.~Gr\"uttm\"uller, R.~S.~Rees and N.~Shalaby,
\emph{Cyclically indecomposable triple systems that are decomposable},
J. Combin. Math. Combin. Comput. \textbf{63} (2007), 103--122.

\bibitem{Lean4}
L.~de~Moura and S.~Ullrich,
\emph{The Lean 4 theorem prover and programming language},
in: Automated Deduction --- CADE 28, Lecture Notes in Comput. Sci.
\textbf{12699}, Springer, 2021, 625--635.

\bibitem{Mathlib}
The mathlib Community,
\emph{The Lean mathematical library},
in: Proc. 9th ACM SIGPLAN Int. Conf. on Certified Programs and Proofs
(CPP 2020), ACM, 2020, 367--381.

\bibitem{MMNS}
M.~Mata-Montero, S.~Normore and N.~Shalaby,
\emph{Generalized Langford sequences: new results and algorithms},
Int. Math. Forum \textbf{9} (2014), 155--181; doi:10.12988/imf.2014.38161.

\bibitem{Nordh}
G.~Nordh,
\emph{Perfect Skolem sets},
Discrete Math. \textbf{308} (2008), 1653--1664; doi:10.1016/j.disc.2006.12.003;
arXiv:math/0506155 (Theorem~15 in the numbering of v1).

\bibitem{OKeefe}
E.~S.~O'Keefe,
\emph{Verification of a conjecture of Th.~Skolem},
Math. Scand. \textbf{9} (1961), 80--82; doi:10.7146/math.scand.a-10624.

\bibitem{Simpson}
J.~E.~Simpson,
\emph{Langford sequences: perfect and hooked},
Discrete Math. \textbf{44} (1983), 97--104; doi:10.1016/0012-365X(83)90008-0.

\bibitem{Skolem}
T.~Skolem,
\emph{On certain distributions of integers in pairs with given differences},
Math. Scand. \textbf{5} (1957), 57--68; doi:10.7146/math.scand.a-10490.

\bibitem{gnlean}
J.~Jones, \emph{gn-lean}: a Lean~4 formalization of the Gurvich--Naumova conjecture (every
$\{1,\dots,n\}$ splits into sets of at most three integers with power-of-three sums), 2026,
whose Langford machinery the present development reuses; registry entry
\texttt{PALOMAR-2026-09-07-000013}, \url{https://github.com/JD-Jones-ASES/gn-lean}.

\bibitem{l2lean}
J.~Jones, \emph{l2-lean}: a Lean~4 formalization of Theorems~A and~B of this note, 2026,
\url{https://github.com/JD-Jones-ASES/l2-lean}.

\end{thebibliography}

\end{document}
